\documentclass[aps,pra,11pt,onecolumn,nofootinbib,tightenlines]{revtex4-2}

\usepackage{mathrsfs}
\usepackage{dsfont}
\usepackage{amsmath}
\usepackage{amsfonts}
\usepackage{amssymb}
\usepackage{amsthm}
\usepackage{xcolor}
\usepackage[colorlinks=true,urlcolor=blue,citecolor=blue,linkcolor=blue]{hyperref}
\usepackage{xurl}
% REVTeX supplies its own caption and longtable* support.  Loading caption
% (directly or through subcaption) redefines longtable* and breaks compilation.
\usepackage{enumerate}
\usepackage{natbib}
\usepackage{bm}
\usepackage{empheq}
\usepackage[shortlabels]{enumitem}

\setlength{\jot}{9pt}
\allowdisplaybreaks

\newtheorem{theorem}{Theorem}[section]
\newtheorem{lemma}[theorem]{Lemma}
\newtheorem{proposition}[theorem]{Proposition}
\newtheorem{corollary}[theorem]{Corollary}
\newtheorem{conjecture}[theorem]{Conjecture}

\theoremstyle{definition}
\newtheorem{definition}[theorem]{Definition}
\newtheorem{assumption}{Assumption}
\newtheorem{example}[theorem]{Example}
\newtheorem{remark}[theorem]{Remark}

\let\oldremark\remark
\renewcommand{\remark}{\oldremark\upshape}

\newcommand{\fo}{\rightarrowtail}
\newcommand{\setFO}{\mathcal{FO}}
\newcommand{\setCO}{\mathcal{CO}}
\newcommand{\setCCO}{\mathcal{CCO}}
\newcommand{\mono}{\mathfrak{R}}
\newcommand{\F}{\mathcal{F}}
\newcommand{\bigket}[1]{\Big| #1 \Big\rangle}
\newcommand{\bigbra}[1]{\Big \langle #1 \Big |}

\newcommand{\mfF}{\mathfrak{F}}
\def \d {\mathrm{d}}
\newcommand{\D}{\mathcal{D}}
\newcommand{\N}{\mathcal{N}}
\newcommand{\M}{\mathcal{M}}
\newcommand{\qstate}{\mathcal{S}(A)}

\renewcommand{\H}{\mathbb{H}}
\newcommand{\cE}{\mathcal{E}}
\newcommand{\bbN}{\mathbb{N}}
\newcommand{\cS}{\mathcal{S}}
\newcommand{\W}{\mathcal{W}}
\newcommand{\X}{\mathcal{X}}
\newcommand{\Prob}{\mathcal{P}}
\newcommand{\V}{\mathcal{V}}
\newcommand{\Y}{\mathcal{Y}}
\newcommand{\cH}{\mathbb{H}}
\newcommand{\cF}{\mathcal{F}}

\newcommand{\id}{I}

\newcommand{\R}{\mathbb{R}}
\newcommand{\T}{\mathbb{T}}


\newcommand{\RR}[1]{\textcolor{red}{(RR: #1)}}

\DeclareMathOperator*{\argmin}{arg\,min}
\DeclareMathOperator*{\argmax}{arg\,max}
\DeclareMathOperator*{\opt}{opt}
\DeclareMathOperator*{\Trm}{Tr}
\DeclareMathOperator*{\Op}{Op \!\!\,}

\def\cE{\mathcal E}

\definecolor{darker}{RGB}{2.0,100.0,20.0}

%\newcommand{\RR}[1]{\textcolor{red}{(RR: #1)}}
\newcommand{\MT}[1]{\textcolor{purple}{(MT: #1)}}
\newcommand{\EH}[1]{\textcolor{darker}{(EH: #1)}}

\renewcommand{\thesection}{\arabic{section}}

%%%%%Erkka's definitions
%fonts
\newcommand{\mc}[1]{\mathcal{#1}}
\newcommand{\ms}[1]{\mathsf{#1}}
\newcommand{\mf}[1]{\mathfrak{#1}}
\newcommand{\mb}[1]{\mathbb{#1}}
%\newcommand{\<}{\langle}
% Semiring preorders
\newcommand{\rleq}{\trianglelefteq}
\newcommand{\rgeq}{\trianglerighteq}


\newenvironment{coloredtext}[1]{\begingroup\color{#1}}{\par\endgroup}

\begin{document}

\title{\Large A complete characterisation of conditional entropies}

\begin{abstract}
Entropies are fundamental measures of uncertainty.  Under natural operational
axioms, ordinary entropy is a barycentre of R\'enyi entropies.  Conditional
entropy asks for the uncertainty seen by an observer with correlated side
information, but its complete axiomatic form has remained open.

We characterize every finite-alphabet conditional entropy that is invariant
under relabeling and embedding, additive on independent systems, normalized on
a fair bit, and monotone under conditionally mixing channels.  Its extremal
members are exponential averages of conditional R\'enyi entropies, together
with their tropical and derivation limits.  We determine the exact admissible
parameter ranges, including the endpoint $\alpha=1$, without imposing a
separate finiteness assumption on the singular weighted moment.

\noindent\textbf{Keywords:} Conditional Entropy, Axiomatic characterization, Resource Theories.

\end{abstract}

\author{Roberto Rubboli$^{*,\dagger,\ddagger}$}
\author{Erkka Haapasalo$^{*}$}
\author{Marco Tomamichel$^{*,\S}$}
\noaffiliation

\maketitle


% Footnotes with matching symbols
\begingroup
\renewcommand{\thefootnote}{\fnsymbol{footnote}}

\footnotetext[3]{Email: \texttt{ror@math.ku.dk}}       % \dagger
\footnotetext[1]{Centre for Quantum Technologies, National University of Singapore, Singapore 117543, Singapore} % *
\footnotetext[2]{Department of Mathematical Sciences, University of Copenhagen, Universitetsparken 5, 2100 Denmark}                         % \ddagger
\footnotetext[4]{Department of Electrical and Computer Engineering,
National University of Singapore, Singapore 117583, Singapore}                         % \S

\endgroup

% A journal submission goes directly to the mathematical content.

\section{Introduction}

For a finite random variable $X$ with probability mass function $P_X$, the
Shannon entropy~\cite{shannon1948mathematical} is
\begin{equation}
 H_1(P_X):=\sum_x P_X(x)\log\frac{1}{P_X(x)}.
 \label{eq:shannon-introduction}
\end{equation}
The logarithm base only fixes the unit; our normalization assigns value one to
a fair bit.  Tail-sensitive uncertainty measures are supplied by the
R\'enyi family~\cite{renyi61}.  For a probability vector $p=(p_i)_{i=1}^d$,
\begin{equation}
 H_\alpha(p):=\frac{1}{1-\alpha}
 \log\!\left(\sum_{i:p_i>0}p_i^\alpha\right),
 \qquad \alpha\in(0,+\infty)\setminus\{1\},
 \label{eq:renyi-introduction}
\end{equation}
with endpoint conventions
\begin{equation}
 H_0(p):=\log|\operatorname{supp}(p)|,\qquad
 H_1(p):=-\sum_{i:p_i>0}p_i\log p_i,\qquad
 H_\infty(p):=-\log\max_i p_i .
 \label{eq:renyi-endpoints-introduction}
\end{equation}
Here $\operatorname{supp}(p):=\{i:p_i>0\}$ and $0\log0:=0$.  The
parameter interval $[0,+\infty]$ is the one-point compactification of the
nonnegative half-line.

Operational axioms single out barycentres of these entropies.  In particular,
every ordinary entropy considered below has the form
\begin{equation}
 \H(P_X)=\int_{[0,+\infty]}H_\alpha(P_X)\,\d\mu(\alpha)
 \label{eq:EntropyBarycentre}
\end{equation}
for a probability measure $\mu$; Section~\ref{sec: set of entropies} gives the
short proof.  Conditional entropy asks for the remaining uncertainty about
$X$ when a correlated register $Y$ is available.  If
$P_{XY}(x,y)=P_Y(y)P_{X|Y=y}(x)$, then
\begin{equation}
 H_1(X|Y)_P=\sum_yP_Y(y)H_1(P_{X|Y=y}).
\end{equation}
Representative conditional R\'enyi entropies include the Hayashi and Arimoto
families and their interpolations
\cite{arimoto1977information,hayashi2011exponential,hayashi2016equivocations}.
Their variety motivates an axiomatic classification rather than the selection
of one formula.

\subsection*{Formal setup and axioms}

All alphabets are finite.  We identify a joint distribution with its
nonnegative matrix
$P_{XY}=(P_{XY}(x,y))_{x\in\X,y\in\Y}$ and write
\begin{equation}
 P_X(x):=\sum_yP_{XY}(x,y),\qquad
 P_Y(y):=\sum_xP_{XY}(x,y),\qquad
 P_{X|Y=y}(x):=\frac{P_{XY}(x,y)}{P_Y(y)}
 \label{eq:joint-marginal-conditional}
\end{equation}
when $P_Y(y)>0$.  A conditional distribution at a zero-mass value of $y$ is
irrelevant and may be chosen arbitrarily.  Relabeling and embedding mean
applying injective maps to the alphabets and assigning zero mass outside their
images.

A mixing channel acts on a probability vector by a doubly stochastic matrix,
namely a nonnegative square matrix whose rows and columns sum to one.  After
embedding two vectors into a common dimension, we write $P_X\succeq Q_X$ when
a mixing channel maps $P_X$ to $Q_X$.  Equivalently, majorization is determined
by the following initial-sum criterion.

\begin{lemma}[Initial sum condition]
\label{lem: Initial sum condition}
Let $P_X,Q_X\in\R^d_{\geq0}$, and let $P_X^\downarrow,Q_X^\downarrow$ denote
their entries in nonincreasing order.  Then $P_X\succeq Q_X$ if and only if
\begin{equation}
 \sum_{i=1}^kP_X^\downarrow(i)\geq
 \sum_{i=1}^kQ_X^\downarrow(i)\qquad(1\leq k<d)
\end{equation}
and the total sums agree.
\end{lemma}

\begin{definition}[Entropy]
\label{def: Entropy}
An entropy is a function
$\H:\bigcup_X\mc P_X\to\R_+$ satisfying:
\begin{enumerate}[itemsep=1pt]
 \item invariance under relabeling and embedding;
 \label{eq:E-2nd-postulate-iso}
 \item monotonicity under mixing:
 $P_X\succeq Q_{X'}\Rightarrow\H(P_X)\leq\H(Q_{X'})$;
 \label{eq:E-1st-postulate-mono}
 \item additivity:
 $\H(P_X\otimes Q_{X'})=\H(P_X)+\H(Q_{X'})$;
 \label{eq:E-2nd-postulate-add}
 \item normalization: $\H((1/2,1/2)^T)=1$.
 \label{eq:E-4-postulate-normalization}
\end{enumerate}
\end{definition}

The conditional operation must allow the mixing applied to $X$ to depend on
side information while permitting arbitrary stochastic processing of $Y$.

\begin{definition}[Conditionally mixing channel and conditional majorization]
\label{def:cond-maj-cond-mix}
A conditionally mixing channel has the matrix form
\begin{equation}
 P_{XY}\longmapsto Q_{XY'}=\sum_{i=1}^kS^{(i)}P_{XY}D^{(i)},
 \label{eq:intro/mixing}
\end{equation}
where every $S^{(i)}$ is doubly stochastic, every $D^{(i)}$ is
sub-stochastic, and $\sum_iD^{(i)}$ is row-stochastic.  We write
$P_{XY}\succeq Q_{X'Y'}$ when, after relabeling and zero-padding both
distributions, such a channel maps the first into the second.
\end{definition}

\begin{definition}[Conditional entropy]
\label{def: conditional entropy}
A conditional entropy is a function
$\H:\bigcup_{X,Y}\mc P_{XY}\to\R_+$ satisfying:
\begin{enumerate}[itemsep=1pt]
 \item invariance under relabeling and embedding;
 \label{eq:QCE-2nd-postulate-iso}
 \item monotonicity under conditional majorization:
 $P_{XY}\succeq Q_{X'Y'}\Rightarrow
 \H(X|Y)_P\leq\H(X'|Y')_Q$;
 \label{eq:QCE-1st-postulate-mono}
 \item additivity:
 $\H(XX'|YY')_{P\otimes Q}
 =\H(X|Y)_P+\H(X'|Y')_Q$;
 \label{eq:QCE-3rd-postulate-add}
 \item normalization:
 $\H(X|Y)_{(1/2,1/2)^T\otimes Q_Y}=1$.
 \label{eq:QCE-4-postulate-normalization}
\end{enumerate}
\end{definition}
We use $\H(P_{XY})$ and $\H(X|Y)_P$ interchangeably.

\subsection*{Main result and proof strategy}

We prove that the extremal conditional entropies are
\begin{equation}
 H_{t,\tau}(X|Y)_P
 =\frac1t\log\sum_{y:P_Y(y)>0}P_Y(y)
 \exp\!\left(t\int_{[0,+\infty]}
 H_\alpha(P_{X|Y=y})\,\d\tau(\alpha)\right),
 \qquad t\neq0,
 \label{eq:General}
\end{equation}
together with the tropical limits $t\to\pm\infty$ and the derivation limit
$t\to0$.  Here $\tau$ is a probability measure on $[0,+\infty]$.  The
finite-temperature ranges are exact.  For $t>0$, monotonicity holds precisely
when
\begin{equation}
 \operatorname{supp}(\tau)\subseteq[0,1],\qquad
 \tau(\{1\})=0,\qquad
 \int_{[0,1)}\frac{\alpha}{1-\alpha}\,\d\tau(\alpha)\leq\frac1t.
 \label{eq:intro-positive-range}
\end{equation}
For $t<0$, either $\operatorname{supp}(\tau)\subseteq[0,1]$, with no
moment-finiteness condition, or there is a single
$\alpha_*\in(1,+\infty]$ such that
$\operatorname{supp}(\tau)\subseteq[0,1]\cup\{\alpha_*\}$,
$\tau(\{1\})=0$, both one-sided integrals
\begin{equation}
 \int_{[0,1)}\frac{\alpha}{1-\alpha}\,\d\tau(\alpha),
 \qquad
 \int_{(1,+\infty]}\frac{\alpha}{1-\alpha}\,\d\tau(\alpha)
 \label{eq:intro-one-sided-moments}
\end{equation}
are finite, and their sum is at most $1/t$.  The two terms are formed
separately, so no undefined $+\infty-(+\infty)$ cancellation is permitted.
Sections~\ref{sec:sufficient-conditions-measure}
and~\ref{sec:necessary-conditions-measures} give the tropical and derivation ranges as
well as the proof sketches for every direction.  Complete analytic proofs are
in Ref.~\cite{rubboli2026parametersFull}; the matching Lean declarations,
dependency map, and statement table are in
Ref.~\cite{rubboli2026parametersLean}.

The proof uses the preordered-semiring framework of
Refs.~\cite{fritz2023abstract,fritz2023abstractII}.  Section~\ref{sec: set of entropies}
first records the unconditional barycentric argument.
Section~\ref{sec: semiring of conditional entropies} constructs the conditional
majorization semiring and determines its homomorphisms and derivations.
Sections~4 and~5 then impose the exact monotonicity ranges: sufficiency follows
from finite-support factorization and common discretization, while necessity
uses two-block, three-block, Shannon-point, and tropical witnesses.
Section~\ref{sec: large sample} converts the resulting spectrum into the
barycentric classification of conditional entropies.

\section{The set of entropies}
\label{sec: set of entropies}
We now prove the barycentric representation~\eqref{eq:EntropyBarycentre} for an
entropy satisfying Definition~\ref{def: Entropy}.  The proof uses the following
large-sample majorization criterion; it is the no-side-information model for
the conditional argument developed later
\cite{mu21_renyi,gour21_axiomatic,haapasalo_inprep}.
\begin{lemma}[{\cite[Proposition 3.7]{Jensen_Kjaerulf_2019}}]\label{lemma:Jensen}
Let $P_X,Q_X$ be two probability distributions and assume that
\begin{equation}
    H_\alpha(P_X) < H_\alpha(Q_X)  \quad \forall \alpha\in[0,+\infty] \,.
\end{equation}
Then, for sufficiently large $n \in \mathbb{N}$ it holds that
$P_X^{\otimes n} \succeq Q_X^{\otimes n}$.
\end{lemma}
This criterion yields a barycentric decomposition against a positive measure,
which normalization makes a probability measure, as stated next.
This result will subsequently serve as a key ingredient in deriving the general
form of a conditional entropy.  Only nonnegative R\'enyi orders occur: negative
orders are incompatible with invariance under embeddings that append zero
coordinates.





\begin{proposition}[General form of entropies]\label{proposition:barycenter entropies}
Let $\H$ be an entropy according to Definition~\ref{def: Entropy}. Then, there exists a probability measure $\mu:\mc B\big([0,+\infty]\big)\to\mb R_+$ such that for all probability distributions $P_X$
\begin{equation}\label{eq:barycenter}
\H(P_X)=\int_{[0,+\infty]}H_\alpha(P_X)\d\mu(\alpha) \,.
\end{equation}
Furthermore, if the normalization axiom~\ref{eq:E-4-postulate-normalization} is not assumed, the conclusion remains valid upon replacing the probability measure with a finite positive measure.
\end{proposition}
\begin{proof}
    The result follows from Lemma \ref{lemma:Jensen}; the proof is essentially the same as that of \cite[Theorem 2]{mu21_renyi} and this result can be seen as a special case of Theorem 7 of \cite{haapasalo_inprep}. We give an outline of this proof for completeness. This same proof also works for Theorem~\ref{th:barycenter conditional entropies} later in this work with slight adaptations.

    The first step is to show that $H_\alpha(P_X)\leq H_\alpha(Q_X)$ for all $\alpha\in[0,\infty]$ implies $\mb H(P_X)\leq\mb H(Q_X)$. First, let us assume that $H_\alpha(P_X)< H_\alpha(Q_X)$ for all $\alpha\in[0,\infty]$. According to Lemma \ref{lemma:Jensen}, this means that, for any $n\in\mb N$ sufficiently large, there is a doubly stochastic matrix $S_n$ such that $S_nP_X^{\otimes n}=Q_X^{\otimes n}$. Thus, since $\mb H$ is monotone and additive,
    \begin{equation}
    n\mb H(P_X)=\mb H(P_X^{\otimes n})\leq\mb H(Q_X^{\otimes n})=n\mb H(Q_X),
    \end{equation}
    i.e., $\mb H(P_X)\leq\mb H(Q_X)$. Let us then assume that $H_\alpha(P_X)\leq H_\alpha(Q_X)$ for all $\alpha\in[0,\infty]$. Denoting by $R$ the two-outcome uniform distribution $(1/2,1/2)$, it follows that
    \begin{equation}
    H_\alpha(P_X^{\otimes m})=mH_\alpha(P_X)\leq mH_\alpha(Q_X)<mH_\alpha(Q_X)+1=H_\alpha(Q_X^{\otimes m}\otimes R\big)
    \end{equation}
    for all $m\in\mb N$ and $\alpha\in[0,\infty]$. Thus, according to our initial step, we also have
    \begin{equation}
    m\mb H(P_X)=\mb H(P_X^{\otimes m})\leq\mb H(Q_X^{\otimes m}\otimes R)=m\mb H(Q_X)+1
    \end{equation}
    for all $m\in\mb N$. This means $\mb H(P_X)\leq\mb H(Q_X)+1/m$ for all $m\in\mb N$, so that $\mb H(P_X)\leq\mb H(Q_X)$, and the initial result is proven.

    For the details of the rest of the proof, we refer to the proof of Theorem 7 of \cite{haapasalo_inprep} and that of Theorem 2 of \cite{mu21_renyi}. The strategy is to define a positive linear functional on the set $C\big([0,\infty]\big)$ of continuous real functions on the compact set $[0,\infty]$ using the initial step proven above and then use the representation result for such functionals.

    Consider the evaluation functions ${\rm ev}_P\in C\big([0,\infty]\big)$, ${\rm ev}_P(\alpha)=H_\alpha(P)$ at any finite probability distribution $P$ and all $\alpha\in[0,\infty]$. By `evaluation' we essentially mean that we actually identify $[0,\infty]$ with the set of R\'{e}nyi entropies $H_\alpha$, $\alpha\in[0,\infty]$. This kind of identification has to be done in the general case of Theorem 7 of \cite{haapasalo_inprep} where the situation is such that large-sample ordering conditions are given by a compact set of functions which is Hausdorff. This is also the case of Theorem \ref{th:barycenter conditional entropies} later on. It can be shown very easily, using our initial step, that we may set up a well-defined positive functional $F$ on the positive cone generated by the evaluation functions in $C\big([0,\infty]\big)$ through $F({\rm ev}_P)=\mb H(P)$. For the well-definedness, note that if $P$ and $Q$ are such that ${\rm ev}_P={\rm ev}_Q$, i.e., $H_\alpha(P)=H_\alpha(Q)$ for all $\alpha\in[0,\infty]$, then, according to our initial result, also $\mb H(P)=\mb H(Q)$. The positivity of $F$ follows similarly. With some extra work, we may show that $F$ extends into a positive functional $G$ on the vector space $V$ spanned by the evaluation functions.

    Suppose that $f\in C\big([0,\infty]\big)$ and fix any $n\in\mb N$ such that $n\geq\|f\|_\infty$. Denoting by $1$ the constant function 1 on $[0,\infty]$ and by $R$ the two-outcome uniform distribution, we have
    \begin{equation}
    f\leq \|f\|_\infty\,1\leq n\,1=n\,{\rm ev}_R={\rm ev}_{R^{\otimes n}},
    \end{equation}
    so that all functions in $C\big([0,\infty]\big)$ are upper bounded by functions from $V$. Now a theorem due to Kantorovich \cite{Kantorovich} related to the Hahn-Banach theorem implies that $G$ extends into a positive linear functional $I:C\big([0,\infty]\big)\to\R$. Due to the Riesz-Markov-Kakutani theorem, there is a finite positive measure $\mu:\mc B\big([0,\infty]\big)\to\R_+$ (the Borel $\sigma$-algebra associated with the one-point compactification of the natural topology of $[0,\infty)$) such that
    \begin{equation}
    I(f)=\int_{[0,\infty]}f\,\mathrm{d}\mu,\qquad\forall f\in C\big([0,\infty]\big).
    \end{equation}
    This means
    \begin{equation}
    \mb H(P)=F({\rm ev}_P)=I({\rm ev}_P)=\int_{[0,\infty]}{\rm ev}_P(\alpha)\,\mathrm{d}\mu(\alpha)=\int_{[0,\infty]}H_\alpha(P)\,\mathrm{d}\mu(\alpha).
    \end{equation}
    Denoting again by $R$ the two-outcome uniform distribution, we have
    \begin{equation}
    1=\mb H(R)=\int_{[0,\infty]}\underbrace{H_\alpha(R)}_{=1}\,\mathrm{d}\mu(\alpha)=\mu\big([0,\infty]\big),
    \end{equation}
    i.e., $\mu$ is a probability measure.
\end{proof}


\section{The semiring of conditional majorization}
\label{sec: semiring of conditional entropies}
In this section, we introduce the main technical tool used to establish sufficient conditions for multiple-copy transformation (large sample result) that lead to the general form of conditional entropies (barycentric decomposition). Our approach relies on the real-algebraic theory of preordered semirings developed in \cite{fritz2023abstract,fritz2023abstractII}. This framework is particularly powerful because it provides a large sample theorem (Theorem \ref{thm:Fritz2022}) whenever one can endow the underlying problem with a suitable semiring structure.
The large sample comparison theorem provides a characterization of the sufficient conditions under which multiple copies of one element of the semiring can be transformed into multiple copies of another.
In our setting, the semiring is essentially the set of joint probability distributions $P_{XY}$, and the large-sample theorem yields necessary and sufficient conditions under which multiple copies of a joint distribution $P_{XY}$ can be transformed into multiple copies of another joint distribution $Q_{XY}$ under conditionally mixing operations in terms of the so-called monotone homomorphisms (see later for the definition) of the theory. In the case of conditional entropies, the monotone homomorphisms and derivations are exactly the extremal conditional entropies (up to a constant and a logarithm). These are the conditional entropies in terms of which any general conditional entropy can be expressed as a convex combination. From here, a standard argument \cite{mu21_renyi,haapasalo_inprep} allows us to deduce from the large sample theorem that any conditional entropy can be expressed as a convex combination of these extremal conditional entropies.

We begin by introducing the fundamental concepts underlying the theory of semirings and the general results that will be employed in our manuscript in Section~\ref{sec: background semiring}. We then construct the semiring associated with conditional entropies and derive the general form of monotone homomorphisms and derivations for this particular setting in Sections~\ref{sec: all monotone homomorphisms}, and~\ref{sec: derivations}.

\subsection{Background on preordered semirings and a \textit{Vergleichsstellensatz}}
\label{sec: background semiring}
We define a preordered semiring as a tuple $S = (\mathcal{S}, +, \cdot, 0, 1, \rleq)$, where $\mathcal{S}$ is a set equipped with binary operations of addition $+$ and multiplication $\cdot$, a zero element $0 \in \mathcal{S}$, a multiplicative unit $1 \in \mathcal{S}$, and a preorder relation $\rleq$  (a reflexive and transitive binary relation) defined on $\mathcal{S}$  satisfying
\begin{equation}
    x \rleq y \ \Rightarrow \
\begin{cases}
x + a \ \rleq\ y + a, \\[0.3em]
x a \ \rleq\ y a,
\end{cases}
\end{equation}
for all $a \in \mathcal{S}$.
Moreover,
$(\mathcal{S}, +, 0)$ and $(\mathcal{S}, \cdot, 1)$ are commutative semigroups, and the multiplication distributes over the addition.
For compactness, we omit the multiplication dot when writing products in the semiring.
We write $x \sim y$ for the equivalence relation induced by $\rleq$, meaning that $x \sim y$ if and only if there exist $z_1, \ldots, z_n \in \mathcal{S}$ such that
\begin{equation}
    x \rleq z_1 \rgeq z_2 \rleq \cdots \rgeq z_n \rleq y.
\end{equation}
A preordered semiring $S$ has polynomial growth if it admits a power universal element $u \in \mathcal{S}$, that is,
\begin{equation}
    x \rleq y \quad \Rightarrow \quad \exists\, k \in \mathbb{N}:\ y \rleq x u^k.
\end{equation}
We say that a preordered semiring $S$ is a preordered semidomain if
\begin{equation}
    \begin{aligned}
&xy = 0 \ \Rightarrow\ x = 0 \ \text{or} \ y = 0, \\
&0 \rleq x \rleq 0 \ \Rightarrow\ x = 0.
\end{aligned}
\end{equation}
Finally, $S$ is zerosumfree if $x + y = 0$ implies $x = 0 = y$.

Given preordered semirings $S$ and $T$, we say that a map $\Phi:\cS\to \mathcal{T}$ is a {\it monotone homomorphism} if it satisfies the properties
\begin{enumerate}
\item Additivity: $\Phi(x+y)=\Phi(x)+\Phi(y)$ for all $x,y\in \cS$,
\item Multiplicativity: $\Phi(xy)=\Phi(x)\Phi(y)$ for all $x,y\in \cS$,
\item Monotonicity: $x\rleq y$ $\Rightarrow$ $\Phi(x)\rleq\Phi(y)$, and
\item $\Phi(0)=0$ and $\Phi(1)=1$.
\end{enumerate}
We call a monotone homomorphism \textit{degenerate} if $ x \rleq y \quad \Rightarrow \quad \Phi(x) = \Phi(y)$. Otherwise, it is called \textit{nondegenerate}.
We need to consider monotone homomorphisms with values in certain special semirings, namely:
\begin{enumerate}
    \item $\mathbb{R}_+$: the half-line $[0,+\infty)$ with the usual addition, multiplication, and total order.
    \item $\mathbb{R}_+^{\mathrm{op}}$: the same set with reversed order.
    \item $\mathbb{T}\mathbb{R}_+$: the half-line $[0,+\infty)$ with the usual multiplication, order, and the tropical sum $x + y := \max\{x,y\}$.
    \item $\mathbb{T}\mathbb{R}_+^{\mathrm{op}}$: the same as for $\mathbb{T}\mathbb{R}_+$, but with reversed order.
\end{enumerate}
The term \emph{temperate reals} refers to the pair $(\mathbb{R}_+, \mathbb{R}_+^{\mathrm{op}})$, while \emph{tropical reals} denotes the pair $(\mathbb{T}\mathbb{R}_+, \mathbb{T}\mathbb{R}_+^{\mathrm{op}})$.
Given a monotone homomorphism $\Phi: \cS \to \mathbb{R}_+$, an additive map $\Delta: \cS \to \mathbb{R}$ is called a derivation at $\Phi$ (or a $\Phi$-derivation) if it satisfies the Leibniz rule
\begin{equation}
    \Delta(xy) = \Delta(x)\,\Phi(y) + \Phi(x)\,\Delta(y)
\end{equation}
for all $x, y \in \cS$.
In this work, we are interested in $\Phi$-derivations at degenerate homomorphisms that are also monotone, i.e.,
\begin{equation}
    x \rleq y \quad \Rightarrow \quad \Delta(x) \leq \Delta(y).
\end{equation}
The following theorem provides both a large-sample and a catalytic result. It belongs to a series of results collectively known as the ``Vergleichsstellens\"{a}tze'' \cite{fritz2023abstractII}.
\begin{theorem}[Based on Theorem 8.6 in \cite{fritz2023abstractII}]\label{thm:Fritz2022}
Let $S$ be a zerosumfree preordered semidomain with a power universal element $u$. Assume that for some $d\in\mb N$ there is a surjective homomorphism $\|\cdot\|:\cS\to\R_{>0}^d\cup\{(0,\ldots,0)\}$ with trivial kernel and such that
\begin{equation}\label{eq:surjectivehomomorphismproperties}
a\rgeq b\ \Rightarrow\ \|a\|=\|b\|\quad {\rm and} \quad \|a\|=\|b\|\ \Rightarrow\ a\sim b.
\end{equation}
Denote the component homomorphisms of $\|\cdot\|$ by $\|\cdot\|_{(j)}$, $j=1,\ldots,d$. Let $x,y\in S\setminus\{0\}$ with $\|x\|=\|y\|$. If
\begin{itemize}
\item[(i)] for every $\mb K\in\{\R_+,\R_+^{\rm op},\T\R_+,\T\R_+^{\rm op}\}$ and every nondegenerate monotone homomorphism $\Phi : \cS \to \mb K$ with trivial kernel, we have $\Phi(x) > \Phi(y)$, and
\item[(ii)] $\Delta(x) > \Delta(y)$ for every monotone $\|\cdot\|_{(j)}$-derivation $\Delta : \cS \to \R$ with $\Delta(u) = 1$ for all component indices $j = 1,\ldots,d$,
\end{itemize}
then
\begin{enumerate}[label=(\alph*),ref=2(\alph*)]
 \item there exists a nonzero $c\in \cS$ such that $cx\rgeq cy$, and
\label{it: catalytic}
\item if additionally $x$ is power universal, then $x^n \rgeq y^n$ for all sufficiently large $n\in\mb N$.
\label{it: asymptotic}
\end{enumerate}
Conversely, if either of these properties holds for at least an integer $n$ or a catalyst $c$, then the inequalities in items (i) and (ii) above hold non-strictly.
\end{theorem}
Large-sample ordering as described in item~\ref{it: asymptotic} of Theorem~\ref{thm:Fritz2022} implies catalytic ordering as in item~\ref{it: catalytic} of Theorem~\ref{thm:Fritz2022}, where the catalyst can be chosen as
\begin{equation}
    c = \sum_{\ell=0}^{n-1} x^\ell y^{n-1-\ell}
\end{equation}
for sufficiently large $n \in \mathbb{N}$. This implication was originally established in \cite{duan2005multiple} for the case where $x, y$ are probability vectors, but the argument extends naturally to the more general framework considered here.
The converse, however, generally fails to hold. In particular,~\cite[Theorem 3]{feng2006relation} presents a method for constructing explicit counterexamples within a specific context.

\subsection{Definition of the conditional majorization semiring}
In this section, we define the semiring of conditional majorization. In particular, we need to specify the elements of the tuple $S=(\mathcal{S}, +, \cdot, 0, 1, \rgeq)$.

\begin{enumerate}

\item \textbf{The set of elements ``$\cS$"}. The elements of the semiring are, essentially, all (not normalized) joint probability distributions $P_{XY}$.
However, since distributions that differ only by embedding into a larger outcome space and/or permutation of outcomes are considered equivalent, the set $\cS$ is formally defined as the set of equivalence classes of joint probability distributions under this similarity relation.
We make this precise in the following.

We can represent (unnormalized) joint distributions as matrices, with rows corresponding to the outcomes of the variable $X$ and columns to those of $Y$. For example, if $X$ has $d$ outcomes and $Y$ has $n$ outcomes, then $P_{XY}$ can be represented as a $d \times n$ matrix
\begin{equation}
P_{XY} = \begin{pmatrix}
P(x_1,y_1) & \hdots & P(x_1,y_n) \\
\vdots &  & \vdots \\
P(x_d,y_1) & \hdots & P(x_d,y_n)
\end{pmatrix} \,.
\end{equation}
We also denote the set of outcomes $\X=\{x_1,...,x_d\}$ and $\Y=\{y_1,...,y_n\}$.
We denote with $\mathcal{V}_{d \times n}$ the set of all $d \times n$ matrices with non-negative entries. Furthermore, define
\begin{equation}
    \mathcal{V} := \bigcup_{d,n \in \mathbb{N}} \mathcal{V}_{d \times n}
\end{equation}
as the set of all matrices with non-negative entries of arbitrary finite dimension.

We say that two joint probability distributions $P_{XY}$ and $Q_{X'Y'}$ are similar, denoted
\begin{equation}
    P_{XY} \approx Q_{X'Y'}  \,,
\end{equation}
if one can be obtained from the other by permuting rows or columns, and/or inserting or deleting rows or columns consisting entirely of zeros.
In other words, $P_{XY}$ and $Q_{X'Y'}$ represent the same distribution up to reordering of outcomes and inclusion or exclusion of impossible events. In yet other words, there is $R_{X''Y''}$ such that both $P_{XY}$ and $Q_{X'Y'}$ can be relabeled and embedded into $R_{X''Y''}$ when we lift the notion of relabeling and embedding to non-normalized distributions.
For $P \in \mathcal{V}$, we define
    $[P] := \{ P' \in \mathcal{V} \mid P \approx P' \}$
as the equivalence class of $P$ under the relation $\approx$. The elements of the semiring are then given by the quotient set
\begin{equation}
    \mathcal{S} := \mathcal{V} / \approx,
\end{equation}
which consists of all equivalence classes of $\mathcal{V}$ under the similarity relation $\approx$.

\item \textbf{The addition ``$+$"}.
Addition in $\mathcal{S} = \mathcal{V}/\!\approx$ is defined via the direct sum of matrices. For two joint distributions $P_{XY}$ and $Q_{X'Y'}$, their direct sum is given by
\begin{equation}
    P_{XY} \oplus Q_{X'Y'} := \begin{pmatrix}
P_{XY} & 0 \\
0 & Q_{X'Y'}
\end{pmatrix}.
\end{equation}
Hence, the sum of two equivalence classes $[P_{XY}]$ and $[Q_{X'Y'}]$ in $\mathcal{S}$ is defined as the equivalence class of the direct sum of any representatives of the respective classes
\begin{equation}
    [P_{XY}] + [Q_{X'Y'}] := [P_{XY} \oplus Q_{X'Y'}] \,.
\end{equation}

\item \textbf{The multiplication ``$\cdot$"}.
Multiplication in $\mathcal{S} = \mathcal{V}/\!\approx$ is defined via the tensor product (or Kronecker product) of two matrices. For $P_{XY} \in \V_{d \times n}$, $Q_{X'Y'} \in \V_{d' \times n'}$
\begin{equation}
P_{XY} \otimes Q_{X'Y'} = \begin{pmatrix}
P(x_1,y_1)Q_{X'Y'} & \hdots & P(x_1,y_n) Q_{X'Y'} \\
\vdots &  & \vdots \\
P(x_d,y_1) Q_{X'Y'} & \hdots & P(x_d,y_n) Q_{X'Y'}
\end{pmatrix} \,,
\end{equation}
where the product of a scalar and a matrix is defined as entry-wise (i.e., element-wise) multiplication. For example,
\begin{equation}
P(x_1,y_1)Q_{X'Y'} = \begin{pmatrix}
P(x_1,y_1)Q(x'_1,y'_1) & \hdots & P(x_1,y_1) Q(x'_1,y'_{n'}) \\
\vdots &  & \vdots \\
P(x_1,y_1) Q(x'_{d'},y'_1) & \hdots & P(x_1,y_1) Q(x'_{d'},y'_{n'})
\end{pmatrix} \,.
\end{equation}
Accordingly, the multiplication of two equivalence classes $[P_{XY}]$ and $[Q_{X'Y'}]$ in $\mathcal{S}$ is defined as the equivalence class of the tensor product of any representatives of the respective classes
\begin{equation}
    [P_{XY}] \cdot [Q_{X'Y'}] := [P_{XY} \otimes Q_{X'Y'}] \,.
\end{equation}

\item \textbf{The zero element ``$0$"}. We define a zero element $0$ in $\V/\!\approx$ as the equivalence class of the $1\times 1$ matrix $0$.

\item \textbf{The unit element ``$1$"}.  We define the unit element $1$ as the equivalence class of the $1\times 1$ matrix $1$.

\item \textbf{The preorder ``$\rgeq$"}.
\label{item: preorder}
We write
\begin{equation}
\label{eq: relation preorder}
\quad[Q_{X'Y'}]\rgeq[P_{XY}].
\end{equation}
if $P_{XY}$ conditionally majorizes $Q_{X'Y'}$, that is, $P_{XY} \succeq Q_{X'Y'}$ in the sense of Definition~\ref{def:cond-maj-cond-mix} when we lift this definition also for non-normalized distributions. Note that in the latter relation, the roles of $P$ and $Q$ are reversed compared to those in equation~\eqref{eq: relation preorder}.
\end{enumerate}

\subsection{Invariance of the homomorphisms under embedding and shifting}
\label{sec: invariance embedding and shifting}
In this section, we prove that monotone homomorphisms are invariant under embeddings of the $X$ register. This invariance is consistent with the requirement that homomorphisms assign a unique value to each element of an equivalent class. Moreover, we show that they are also invariant under column-wise shifts.
This property will be crucial in characterizing the general form of such homomorphisms, which in turn determines the general form of conditional entropies.

\bigskip
We begin by introducing some notation. Let $0_{n \times 1}$ denote the column vector of zeros of dimension $n$, $0_{1 \times n}$ the corresponding row vector, and $0_{n \times n}$ the $n \times n$ zero matrix. We denote by $E^{j \times j}_{n \times n}$ the matrix of size $n \times n$ with a 1 in the $(j, j)$-th entry and zeros elsewhere. Finally, $I_{n \times n}$ stands for the $n \times n$ identity matrix.


The next lemma states that the columns can be shifted without affecting the value of any homomorphism, as there exists a channel implementing the transformation in both directions.
\begin{lemma}[Shifting]\label{lemma:shifting}
    Let $P_{XY} \in \V_{d\times n}$. Moreover,  for each $i=1,...,n$, let $P_{X,Y=y_i}$ be the $i$-th column of $P_{XY}$, that is, $P_{XY} = (P_{X,Y=y_1},...,P_{X,Y=y_n})$. Define
    \begin{equation}
\widetilde{P}_{\widetilde{X}Y}=\bigoplus_{i=1}^n P_{X,Y=y_i}= \begin{pmatrix}
P_{X,Y=y_1} & 0_{d \times 1}  & \hdots & 0_{d \times 1} \\
0_{d \times 1} & P_{X,Y=y_2} & \hdots & 0_{d \times 1} \\
\vdots &  & \vdots & \vdots\\
0_{d \times 1} & \hdots & 0_{d \times 1} & P_{X,Y=y_n}\,
\end{pmatrix} \,,
\end{equation}
where $\widetilde{X}=X\sqcup\cdots\sqcup X$ is the $n$-fold separate union of $X$. We have
\begin{equation}
 \begin{pmatrix}
        P_{XY}\\
        0_{dn \times n}
    \end{pmatrix}\succeq\widetilde{P}_{\widetilde{X}Y}\succeq  \begin{pmatrix}
        P_{XY}\\
        0_{dn \times n}
    \end{pmatrix}.
\end{equation}
\end{lemma}

\begin{proof}
Let us first prove that
\begin{equation}
    \begin{pmatrix}
        P_{XY}\\
        0_{dn \times n}
    \end{pmatrix}\succeq\widetilde{P}_{\widetilde{X}Y}\,.
\end{equation}
For this case, the proof follows by first measuring the $Y$ register and, depending on the outcome, performing an appropriate relabeling of the outcomes in the $X$ register.
 Explicitly, let us define $D^{(j)}= E^{j \times j}_{n \times n}$ which selects the outcome $y_j$ of the $Y$ register. For example, for $j=2$, we have
\begin{equation}
    D^{(2)}=\begin{pmatrix}
        0 & 0 & 0_{1 \times (n-2)}  \\
       0 & 1 & 0_{1 \times (n-2)} \\
       0_{(n-2) \times 1} & 0_{(n-2) \times 1} & 0_{(n-2) \times (n-2)}
    \end{pmatrix} \,.
\end{equation}
The matrix $S^{(j)}$ relabels the outcomes on the $X$ register. For example, for $j=2$, we have
\begin{equation}
    S^{(2)}=\begin{pmatrix}
        0_{d \times d} & I_{d \times d} & 0_{d \times d(n-2)} \\
        I_{d \times d} & 0_{d \times d} & 0_{d \times d(n-2)} \\
        0_{d(n-2) \times d} & 0_{d(n-2) \times d} & I_{d(n-2) \times d(n-2)}
    \end{pmatrix} \,.
\end{equation}
To prove the other direction, we use the same matrices as before to select the appropriate outcome of $Y$ and then apply the same permutation to restore the original ordering of the entries.
\end{proof}

\begin{remark}\label{rem:SomOfColumns}
Lemma \ref{lemma:shifting} essentially tells us that, for a matrix $P_{XY}\in\V_{d \times n}$ with columns $P_{X,Y=y_i}$, we may write
\begin{equation}
[P_{XY}]\rgeq[P_{X,Y=y_1}]+\cdots+[P_{X,Y=y_n}]\rgeq [P_{XY}] \,.
\end{equation}
This will subsequently aid us in finding all the monotone homomorphisms of $S$ into $\mb K\in\{\mb R_+,\mb R_+^{\rm op},\mb{TR}_+,\mb{TR}_+^{\rm op}\}$.
\end{remark}

\subsection{Some important properties of the semiring \texorpdfstring{$S$}{S}}
To use the large-sample results in Theorem~\ref{thm:Fritz2022}, we must first show that the semiring is of polynomial growth, i.e., it contains a power-universal element. Since the large sample results apply only when the input element is power-universal, we also need to determine exactly which elements have this property. Lastly, we must construct a surjective degenerate homomorphism $\|\cdot\|:\cS\to\R_{>0}^d\cup\{0\}$ with the properties required in Theorem~\ref{thm:Fritz2022}. We begin with the last point. In the following lemma, we show that such a homomorphism can be chosen as the total weight, given by the sum of the matrix entries associated with elements of the semiring.
Finally, note that the total weight is a scalar quantity rather than a vector, and therefore there exists only a single component homomorphism of \(\|\cdot\|\) in Theorem~\ref{thm:Fritz2022}, namely \(j = 1\).
\begin{lemma}
\label{lem: degenerate and zig-zag}
In the preordered semiring of conditional majorization $S$, the map
\begin{equation}
\|[P_{XY}]\| = \sum_{x\in\mc X}\sum_{y\in\mc Y} P_{XY}(x,y) \,,
\end{equation}
where $P_{XY}$ can be chosen to be any representative of the equivalence class $[P_{XY}]$,
defines a surjective homomorphism $\|\cdot\| : \cS \mapsto \mathbb{R}_+$ with trivial kernel. Moreover, it satisfies
\begin{equation}
[P] \rgeq [Q] \quad \Rightarrow \quad \|[P]\| = \|[Q]\| \quad \Rightarrow\quad  [P] \sim [Q]\,.
\end{equation}
\end{lemma}
\begin{proof}
Clearly, the total weight is a homomorphism since $\|[P]\cdot [Q]\|=\|[P\otimes Q]\|=\|[P]\|\|[Q]\|$ and $\|[P]+[Q]\|=\|[P\oplus Q]\|=\|[P]\|+\|[Q]\|$ for all $[P],[Q]\in\cS$.

The first direction, namely that $[P] \rgeq [Q] \Rightarrow \|[P]\| = \|[Q]\|$, follows immediately, since conditionally mixing operations preserve the total weight $\|P\|$ of any element $P$. This also shows that the total weight is a degenerate homomorphism since it is invariant under the preorder $\rgeq$.

Let us now turn to the implication $\|[P]\|=\|[Q]\|\Rightarrow [P]\sim [Q]$.  We show that if we define $\|P\|=\|Q\| = c$, then $P\preceq c\succeq Q$
and hence $[P]\sim [Q]$. In particular, we show that from the element $[c]$ we can generate both $P$ and $Q$. This is straightforward by considering the matrix with a single non-zero entry $c$ in the top-left corner. Indeed, starting from this matrix, we can construct both $P$ and $Q$. For example, let us focus on $P_{XY} \in \V_{d\times n}$.
Then the steps are the following. Let us denote with $P_{Y}$ the row vectors with components $P_Y(y_i) = \sum_jP_{XY}(x_j,y_i)$.
First, we create for each outcome $Y$ a column whose first entry is $P_Y(y_i)$ and all other entries are zero. This is done by multiplying on the r.h.s by a row stochastic matrix.
\begin{align}
\begin{pmatrix}
     P_{Y} \\
    0_{(d-1) \times n}
\end{pmatrix}=
\begin{pmatrix}
    \|P_{XY}\| & 0_{1\times (n-1)} \\
    0_{(d-1)\times 1} & 0_{(d-1)\times (n-1)}
\end{pmatrix}
\begin{pmatrix}
    \|P_{XY}\|^{-1} P_{Y} \\
    \vdots \\
    \|P_{XY}\|^{-1} P_{Y}
\end{pmatrix}\,.
\end{align}
Next, we measure the register $Y$ and, conditioned on the outcome $y_i$, we transform the deterministic column (with a single non-zero entry) into a more distributed column $P_{X,Y=y_i}$. Let us denote with $P_{X,Y=y_1}$ the column with components $P_{X,Y=y_1}(x_j)=P_{X,Y}(x_j,y_1)$ and with $P_{X|Y=y_1}$ the column with components $P_{X|Y=y_1}(x_j)=P_{X,Y}(x_j,y_1)/P_Y(y_1)$,
For the first outcome $y_1$, we obtain
\begin{align}
\begin{pmatrix}
     P_{X,Y=y_1} & 0_{d \times (n-1)}
\end{pmatrix}=
\begin{pmatrix}
    P_{X|Y=y_1} & \star_{1\times (d-1)}
\end{pmatrix}
\begin{pmatrix}
     P_{Y} \\
    0_{(d-1) \times n}
\end{pmatrix}
\begin{pmatrix}
    1 & 0_{1\times (n-1)}\\
    0_{(n-1) \times 1} & 0_{1\times (n-1)}
\end{pmatrix}
\end{align}
Here, $\star_{1\times (d-1)}$ is any entry that completes the matrix to a doubly stochastic one. By summing over all the outcomes $y_i$ for all $i=1,..,n$, we obtain the desired probability $P_{XY}$. This shows that overall, we have
\begin{equation}
    P_{XY} = \sum_{i=1}^n S^{(i)} P DD^{(i)} \,,
\end{equation}
where
\begin{align}
P=\begin{pmatrix}
    \|P_{XY}\| & 0_{1\times (n-1)} \\
    0_{(d-1)\times 1} & 0_{(d-1)\times (n-1)}
\end{pmatrix} \,, \quad
    D=\begin{pmatrix}
    \|P_{XY}\|^{-1} P_{Y} \\
    \vdots \\
    \|P_{XY}\|^{-1} P_{Y}
\end{pmatrix} \,,
\end{align}
and for example, for $i=1$, we have
\begin{align}
D^{(1)}=\begin{pmatrix}
    1 & 0_{1\times (n-1)}\\
    0_{(n-1) \times 1} & 0_{1\times (n-1)}
\end{pmatrix} \,,\quad
S^{(1)} = \begin{pmatrix}
    P_{X|Y=y_1} & \star_{1\times (d-1)}
\end{pmatrix}  \,.
\end{align}

\end{proof}

As we mentioned above, to apply the large sample and catalytic results of Theorem~\ref{thm:Fritz2022}, we must also identify a power universal element. By definition, this implies that the semiring is of polynomial growth. As defined in Section~\ref{sec: background semiring}, a power universal element is one that, when provided in sufficient copies, enables the realization of any transformation between elements of the semiring. Within the semiring of conditional majorization, such an element corresponds to any probability distribution exhibiting non-zero randomness in the $X$ register—that is, a probability distribution that is not fully deterministic in $X$. Furthermore, this condition must be satisfied for every possible outcome of the $Y$ register.

\begin{proposition}
\label{prop: power universals}
The following statements are equivalent for $[R_{XY}] \in \cS$:
\begin{enumerate}
    \item \( [R_{XY}] \in \cS \) is power universal.
    \label{item: power universal}
     \item  \( \|[R_{XY}]\| = 1 \), and every non-zero column \( R_{X,Y=y} \) of any representative \( R_{XY} \) contains at least two non-zero entries.
    \label{item: two non-zero entries}
    \item There exists \( r \in (0,1) \) such that
    \begin{equation}\label{eq:pucond}
    \begin{pmatrix}
        r \\
        1 - r
    \end{pmatrix}
    \succeq R_{XY}.
    \end{equation}
    \label{item: majorization R_r}
\end{enumerate}
Consequently, the semiring $S$ is of polynomial growth.
\end{proposition}

\begin{proof}
We prove separately the implications $\ref{item: power universal} \implies \ref{item: two non-zero entries}$, $\ref{item: two non-zero entries} \implies \ref{item: majorization R_r}$ and $\ref{item: majorization R_r} \implies \ref{item: power universal}$. This shows that the statements are all equivalent.
In the following, for brevity, we simply write $P_{XY}$ to refer to a representative of the equivalence class $[P_{XY}]$, without always stating it explicitly. Throughout this proof, we assume that $R_{XY}\in\mc V_{d\times n}$ (i.e.\ $|\mc X|=d$ and $|\mc Y|=n$) and we number $\mc X$ and $\mc Y$ explicitly to aid our matrix calculations.

\bigskip

Let us start with the implication $\ref{item: power universal} \implies \ref{item: two non-zero entries}$. The fact that \( \|[R_{XY}]\| = 1 \) is immediate, since conditional majorization preserves the total weight of elements. If instead \( \|[R_{XY}]\| \neq 1 \), the transformation would not be possible. Next, we show that if $R_{XY}$ contains a column with only one non-zero entry, then it cannot be a power universal. The key idea is that if \( R_{XY} \) contains a column with only one non-zero entry, then any tensor power \( R_{XY}^{\otimes m} \) will also contain a column of the same form. However, starting from a mixed probability distribution, it is impossible to generate a deterministic distribution conditioned on the event that \( Y \) corresponds to such a column. This follows from the fact that mixing channels cannot produce a deterministic distribution from a non-deterministic one; instead, they only increase the level of mixing.
Explicitly, let us define \( P_{XY} = (1/2, 1/2) \) and take \( Q_{XY} = 1 \) to be the single unit element. Note that we use the subscript $XY$ for both, as the smaller matrix can always be embedded into the space of the larger one by padding with zeros.
Observe that \( Q_{XY} \succeq P_{XY} \).  Assume now that \( R_{XY} \) is power universal. Then, there exists a sufficiently large \( m \in \mathbb{N} \) such that
\begin{equation}
    P_{XY} \succeq Q_{XY} \otimes R_{XY}^{\otimes m} = R_{XY}^{\otimes m}\,.
\end{equation}
This implies that, for a matrix \( R_{XY} \in \V_{d \times n} \), there exist an integer \( k \in \mathbb{N} \), doubly stochastic matrices \( S^{(1)}, \ldots, S^{(k)}\), and (without loss of generality—since the zero entries in the embedding will eliminate the effect of the remaining components) single-row matrices \( D^{(i)} = (d^{(i)}_1, \ldots, d^{(i)}_{n^m}) \),
such that an embedding of the vector \( P_{XY} \) can be transformed into \( R_{XY}^{\otimes m} \)
\begin{align}\label{eq:apu3}
R_{XY}^{\otimes m}&=\sum_{i=1}^k S^{(i)}
\begin{pmatrix}
    \frac{1}{2}\\
    \frac{1}{2} \\
    0_{(d^m-2)\times 1}
\end{pmatrix}
\big(d^{(i)}_{1}\ \cdots\ d^{(i)}_{n^m}\big)=\sum_{i=1}^k S^{(i)}
\begin{pmatrix}
    \frac{1}{2}d^{(i)}_{1}\ \cdots\ \frac{1}{2}d^{(i)}_{n^m} \\
   \frac{1}{2} d^{(i)}_{1}\ \cdots\ \frac{1}{2}d^{(i)}_{n^m} \\
   0_{(d^m-2) \times 1}
\end{pmatrix} \,.
\end{align}
We now distinguish two cases: either \( d_j^{(i)} = 0 \) or \( d_j^{(i)} > 0 \). In both situations, the column of \( R_{XY}^{\otimes m} \) with a single non-zero entry cannot be generated by any column of the matrix \( P_{XY} D^{(i)} \). Indeed, if \( d_j^{(i)} = 0 \), the corresponding column is identically zero; if \( d_j^{(i)} > 0 \), then the resulting column contains at least two non-zero entries of the form \( \frac{1}{2} d_j^{(i)} \), and a doubly stochastic matrix cannot reduce this uncertainty to produce a deterministic outcome. This follows from the fact that both the rows and columns of a doubly stochastic matrix must sum to one, thus preserving entropy rather than eliminating it.

\bigskip
Let us now turn to the implication $\ref{item: two non-zero entries} \implies \ref{item: majorization R_r}$
Let us first assume that $\|R_{XY}\|=1$ and $R_{XY}$ is such that its columns have supports with at least two points. We need to find $r\in(0,1)$ such that the condition in item (3) holds.
The key idea is that each column of \( R_{XY} \) containing at least two non-zero entries can be generated from a vector of the form \( (r, 1 - r) \), for some sufficiently large \( r \in (0,1) \). Such a vector represents an almost deterministic distribution, which—after suitable embedding—can be transformed, via mixing, into any non-deterministic distribution with the same support. Furthermore, the same vector \( (r, 1 - r) \) can be used to generate all columns of \( R_{XY} \), since we are free to enlarge the system \( Y \) to accommodate additional outcomes if needed. Explicitly,
because of the restriction on $R_{XY}$, we have $R_{XY}(x_j,y_i)<R_Y(y_i)$ for all $x_j\in \mathcal{X}$ and $y_i\in \mathcal{Y}$ such that the column $R_{X,Y=y_i}$ is non-zero. Therefore, there exists a $r\in[1/2,1)$ such that
\begin{equation}
R_Y(y_i)r\geq R_{XY}(x_j,y_i)\qquad\forall(x_j,y_i)\in \mathcal{X}\times \mathcal{Y}:\ R_{X,Y=y_i}\neq0.
\end{equation}
By Lemma~\ref{lem: Initial sum condition}, it follows that for each \( y_i \in \mathcal{Y} \) with \( R_{X,Y=y_i} \neq 0 \), the two-dimensional embedded vector $\widetilde{R}_Y(y_i) := \big(r R_Y(y_i), (1 - r) R_Y(y_i)\big)$
majorizes the column \( R_{X,Y=y_i} \).
Hence, for each \( y_i \in \mathcal{Y} \) with \( R_{X,Y=y_i} \neq 0 \), there exists a doubly stochastic matrix \( S^{(i)} \) and an embedding of the vector \( \widetilde{R}_Y(y_i) \), denoted by \( \widehat{R}_Y(y_i) \), such that $S^{(i)} \widehat{R}_Y(y_i) = R_{X,Y=y_i}$.
To isolate the outcome \( y_i \in \mathcal{Y} \), we define the matrices \( D^{(i)} \) with a $1$ in the \( (y_i, y_i) \)-th position and zeros elsewhere. Treating \( R_Y \) as a row-stochastic vector and indexing the outcomes of $Y$ as \( i = 1,..,n \), we obtain
\begin{align}
\left(\begin{array}{c}
r\\1-r
\end{array}\right)&\succeq\left(\begin{array}{c}
r\\1-r
\end{array}\right)R_Y\\
&=\left(\begin{array}{ccc}
rR_Y(y_1)&\cdots&rR_Y(y_n)\\
(1-r)R_Y(y_1)&\cdots&(1-r)R_Y(y_n)
\end{array}\right)\\
&\succeq\sum_{i=1}^n S^{(i)}\left(\begin{array}{ccc}
\widehat{R}_Y(y_1)\cdots \widehat{R}_Y(y_n)
\end{array}\right)D^{(i)}\\
&=\big(S^{(y_1)}\widehat{R}_Y(y_1)\cdots S^{(y_n)}\widehat{R}_Y(y_n)\big)\\
&=\big(R_{X,Y=y_1}\cdots R_{X,Y=y_n}\big)\\
&=R_{XY}.
\end{align}
Hence, we obtain what we wanted to prove.

\bigskip

Finally, let us prove the implication $\ref{item: majorization R_r} \implies \ref{item: power universal}$. We need to show that if there exists $r\in(0,1)$ such that $R_r\succeq R_{XY}$, where $R_r$ is the matrix (column) appearing on the left-hand side of \eqref{eq:pucond},
then $[R_{XY}]$ is power universal. To establish this, we first show that $[R_r]$ is a power universal element; that is, for any elements $[P_{XY}]$ and $[Q_{XY}]$ such that $\|[P_{XY}]\| = \|[Q_{XY}]\|$, there exists a sufficiently large integer $m \in \mathbb{N}$ such that
$P_{XY} \succeq Q_{XY} \otimes R_r^{\otimes m}
$. This, in turn, implies that [$R_{XY}]$ is also power universal, since $R_r\succeq R_{XY}$ implies the chain of inequalities
\begin{equation}
    P_{XY} \succeq Q_{XY} \otimes R_r^{\otimes m}\succeq Q_{XY} \otimes R_{XY}^{\otimes m}\,.
\end{equation}
Note that we use the subscript $XY$ for both $P_{XY}$ and $Q_{XY}$, since the smaller matrix can always be embedded into the larger one by zero-padding.
Without loss of generality, we may assume that $r \in [1/2, 1)$, as the case $r \in (0, 1/2)$ lies in the same equivalence class due to the ability to permute elements.
The fact that $[R_r]$ is power universal follows from the observation that, given a sufficiently large number of copies of $R_r^{\otimes m}$, the non-zero entries of the resulting vector become arbitrarily small. In particular, the largest entry is $r^m$, which approaches zero as $m \to \infty$.
By the initial sum condition for majorization (see Lemma~\ref{lem: Initial sum condition}), this implies that any column of a suitably large embedded distribution $P_{X,Y=y_i}$ can be transformed into $Q_{X,Y=y_i} \otimes R_r^{\otimes m}$, provided the corresponding weights match; that is, when $P_Y(y_i) = Q_Y(y_i)$.
If the column weights differ, we can apply an additional stochastic map to $P_{XY}$ to adjust the marginals accordingly. Explicitly, let $D$ be a row-stochastic matrix such that
\begin{equation}
    P_{Y}D=Q_Y
\end{equation}
and denote $\widetilde{Q}_{XY}:=P_{XY}D$. It is easy to notice that the columns $\widetilde{Q}_{X,Y=y}$ have the same weight as the columns $Q_{X,Y=y}$ for each $y$.
Moreover,  let $m\in\mb N$ be large enough so that
\begin{equation}
r^m\max_{x \in\mc X}Q_{XY}(x,y)\leq \min_{\substack{x \in\mc X \\ \widetilde{Q}(x, y) \neq 0}} \widetilde{Q}(x, y) \qquad\forall y\in\mc Y.
\end{equation}
It then follows, by Lemma~\ref{lem: Initial sum condition}, that for every $y \in\mc Y$, the column $\widetilde{Q}_{X,Y=y}$ majorizes the column $Q_{X,Y=y} \otimes R^{\otimes m}$.
Let $S^{(i)}$ be a doubly stochastic matrix such that
\begin{equation}
    S^{(i)} \widehat{Q}_{X,Y=y_i} = Q_{X,Y=y_i} \otimes R^{\otimes m} \,,
\end{equation}
where we denoted with $\widehat{Q}_{X,Y}$ a suitably embedding of $\widetilde{Q}_{XY}$.
Furthermore, let $D^{(i)}$ denote the matrix with a 1 at entry $(y_i, y_i)$ and zeros elsewhere.
We now have
\begin{align}
P_{XY}&\succeq P_{XY}D=\widetilde{Q}_{XY}\succeq\sum_{i=1}^kS^{(i)}\widehat{Q}_{XY}D^{(i)} =  Q_{XY}\otimes R^{\otimes m}\,.
\end{align}
Finally, the semiring has polynomial growth, as it evidently contains at least one power-universal element.

\end{proof}
Let us note that, whenever $Q_Y$ is a probability distribution, the joint distribution $(1/2,1/2)\times Q_Y$ corresponds to a power universal of the semiring $S$ according to the above result. Note that this is exactly the kind of joint distribution that we use to normalize our general conditional entropies in the normalization axiom 4 of Definition~\ref{def: conditional entropy}.

\subsection{The homomorphisms}
\label{sec: all monotone homomorphisms}
In this section, we derive the general form of conditional entropy satisfying the axioms stated in Definition~\ref{def: conditional entropy}. Within the theory of preordered semirings, monotone homomorphisms and derivations are functions designed to satisfy analogous properties. In particular, for the semiring of conditional majorization, these functions correspond—up to an additive constant and a logarithmic factor—to the extremal conditional entropies, which we later denote by $H_{t,\tau}$. As we prove later, these entropies generate the entire family of conditional entropies via convex combinations (see Theorem~\ref{th:barycenter conditional entropies}). Hence, for our purposes, it is crucial to determine the form of the homomorphisms and derivations, which is achieved by characterizing them through their mathematical properties.



As we reviewed in Section~\ref{sec: background semiring}, we need to consider monotone homomorphisms \(\Phi: S \to \mathbb{K}\) in the fields $
    \mathbb{K} = \mathbb{R}_+,\mathbb{R}_+^{\mathrm{op}},\mathbb{TR}_+,\mathbb{TR}_+^{\mathrm{op}} \,.$
The cases $\mathbb{R}_+,\mathbb{R}_+^{\mathrm{op}}$ correspond to the temperate case, whereas $\mathbb{TR}_+,\mathbb{TR}_+^{\mathrm{op}} $ correspond to the tropical cases.
As we show later, the temperate cases correspond to the conditional entropy $H_{t,\tau}$ (up to a constant and a logarithmic factor), with the probability measure $\tau$ and $t>0$ on $\mathbb{R}_+$ and $t<0$ on $\mathbb{R}_+^{\mathrm{op}}$ (see Definition~\ref{def: extremal entropies}). The tropical cases yield the limits $t\to+\infty,-\infty$. In particular, $\mathbb{K} = \mathbb{TR}_+^{\rm op}$ gives rise to the entropies $H_{-\infty,\tau}$ (see Definition~\ref{def: extremal entropies}), corresponding to the limiting form of $H_{t,\tau}$ as $t\to-\infty$.
Finally, the case \(\mathbb{TR}_+\) gives rise to \(H_{+\infty,0}\), which corresponds to the limiting form of \(H_{t,\tau}\) as \(t \to +\infty\) and to the choice of \(\tau\) as a point measure at \(\alpha = 0\) (see Definition~\ref{def: extremal entropies}). As discussed later, in this last limiting regime, the convexity properties of the function enforce that only point measures supported at \(\alpha = 0\) are admissible. In the following result, we express the statement in terms of a measure \(\mu\), which we later decompose as \(\mu = t\,\tau\), where \(t\) denotes the total mass of \(\mu\) and \(\tau\) is the associated probability measure.


\begin{proposition}
\label{prop: monotone homomorphisms}
Let \(S\) be the semiring of conditional majorization. The nondegenerate
monotone homomorphisms with trivial kernel
$\Phi:\mc S\to\mb K\in\{\R_+,\R_+^{\rm op},\T\R_+,\T\R_+^{\rm op}\}$ are of the form
\begin{equation}
\label{eq:FinalForm}
\Phi(P_{XY}) =
\begin{cases}
\displaystyle \sum_{y \in \mathcal{Y}} P_Y(y) \exp\Bigg(\int_{[0,+\infty]} H_\alpha(P_{X|Y=y}) \, \mathrm{d}\mu(\alpha) \Bigg), & \mb K\in\{\mathbb{R}_+,\mathbb{R}_+^{\mathrm{op}}\}, \\
\displaystyle \max \limits_{\substack{y \in \mathcal{Y} \\ P_Y( y) > 0}} \exp\Bigg(\int_{[0,+\infty]} H_\alpha(P_{X|Y=y}) \, \mathrm{d}\mu(\alpha) \Bigg), &\mb K=\mathbb{TR}_+^{\mathrm{op}}, \\
\displaystyle \max \limits_{\substack{y \in \mathcal{Y} \\ P_Y( y) > 0}} \exp\left(k H_0(P_{X|Y=y}) \, \right), &\mb K=\mathbb{TR}_+ \,,
\end{cases}
\end{equation}
where $\mu$ is a finite positive measure in case $\mb K=\mathbb{R}_+$, and it is a negative measure in cases $\mb K\in\{\mathbb{R}_+^{\mathrm{op}},\T\R_+^{\rm op}\}$. Moreover, in case $\mb K=\mathbb{TR}_+$,  $k$ is a positive constant.
\end{proposition}


\begin{proof}
The proof is organized into several key steps that lead to the general form of the monotone homomorphisms. Specifically, we proceed as follows:
\begin{enumerate}
    \item Step 1: Additivity. We use the shifting lemma~\ref{lemma:shifting} to reduce the problem to determining the monotone homomorphism for a fixed column vector \( P_{X,Y=y} \), corresponding to a specific value \( Y = y \).
    \item Step 2: Multiplicativity. We then use the property that monotone homomorphisms must be multiplicative under tensor products to further simplify the problem to normalized conditional distributions \( P_{X|Y=y} \).
    \item Step 3: Inner barycenter. We invoke Proposition~\ref{proposition:barycenter entropies}, which states that the most general form of an entropy can be expressed as a linear combination of Rényi entropies.
     \item Step 4: Simplifying $\mathbb{TR}_+$. We use the fact that, in this case, the monotone homomorphisms are strongly quasi-concave, which allows us to further simplify their form and to single out point measures supported at $\alpha=0$.
\end{enumerate}

\bigskip
\bigskip


\noindent\textbf{Step 1: Additivity}. The first step is to use the additivity of the homomorphism under direct sums to reduce the problem from the matrix $P_{XY}$ to a single column vector $P_{X,Y=y}$ for a fixed value $Y = y$.
Lemma~\ref{lemma:shifting}, together with the fact that the homomorphism is additive in both the temperate and tropical sense under summation (see also Remark~\ref{rem:SomOfColumns}), leads to the following additive form
\begin{equation}\label{eq:additivity}
\Phi(P_{XY}) =
\begin{cases}
\sum_{y \in \Y} \Phi(P_{X,Y=y}), & \mathbb{R}_+,\mathbb{R}_+^{\mathrm{op}}, \\
\max_{y \in \Y} \Phi(P_{X,Y=y}), & \mathbb{TR}_+,\mathbb{TR}_+^{\mathrm{op}} \,.
\end{cases}
\end{equation}
Hence, the problem reduces to finding the homomorphisms for the column vector $P_{X,Y=y}$ with fixed $Y=y$.

\bigskip
\bigskip
 \noindent\textbf{Step 2: Multiplicativity}. The second step applies the multiplicativity of the homomorphism under tensor products, specifically scalar multiplication, to reduce the problem from the unnormalized vector $P_{X,Y=y}$ to the normalized vector probability distribution $P_{X|Y=y} = P_{X,Y=y}/P_Y(y)$, where $P_Y(y) = \sum_x P_{XY}(x,y) = P_Y(y)$ is the total weight of the column $P_{X,Y=y}$. We then obtain
 \begin{equation}\label{eq:normalize}
     \Phi(P_{X,Y=y}) = \Phi\left(P_Y(y) \otimes \frac{P_{X,Y=y}}{P_Y(y)}\right) =\Phi(P_Y(y)) \Phi\left(\frac{P_{X,Y=y}}{P_Y(y)}\right) = \Phi(P_Y(y)) \Phi\big(P_{X|Y=y}\big) \,.
 \end{equation}

Next, we analyze the form of $\Phi(P_Y(y))$, which corresponds to the homomorphism applied to the single matrix entry $P_Y(y) \in \mb R_+$.
By definition, the homomorphism satisfies $\Phi(0) = 0$ for the zero element and $\Phi(1) = 1$ for the unit element.
Moreover, from the multiplicativity of $\Phi$, it follows that $\Phi(xy)=\Phi(x)\Phi(y)$ for all $x,y\in\mb R_+$. Thus, if the restriction of $\Phi$ to $\mathbb{R}_+$ satisfies a suitable regularity condition—such as being bounded on an interval with non-empty interior—it follows that $\Phi(x) = x^\gamma$ for some $\gamma \in \mathbb{R}$.
 We shall next show that $\Phi(x)\leq 1$ for all $x\in[0,1]$. Let $x\in[0,1]$. We easily see that
\begin{equation}\label{eq:x1}
\left(\begin{array}{cc}
x&0\\
0&1-x
\end{array}\right)\succeq\left(\begin{array}{cc}
1&0\\
0&0
\end{array}\right)\succeq\left(\begin{array}{cc}
x&0\\
0&1-x
\end{array}\right) \,.
\end{equation}
The first implication follows from the fact that we can apply a mixing channel to the second column, conditioned on the second outcome of the register $Y$, and then merge both outcomes into a single $Y$ outcome. The second implication follows by applying a stochastic map on $Y$ that discards the second outcome and creates two new outcomes with probabilities $x$ and $1-x$. Finally, we permute the outcomes of the $X$ register conditioned on the second outcome of the $Y$ register. Since the homomorphism must be monotone under both actions, equality must hold for both elements, and thus, for the temperate case, the sum property implies that
\begin{equation}
1=\Phi(1)=\Phi\left[\left(\begin{array}{cc}
x&0\\
0&1-x
\end{array}\right)\right]=\Phi(x)+\Phi(1-x)\geq\Phi(x) \,,
\end{equation}
since $\Phi(1-x)\geq 0$. The tropical case is similar since $\max\{\Phi(x),\Phi(1-x)\} \geq \Phi(x)$. Hence, the homomorphisms are bounded and, for what was discussed earlier, $\Phi(x)=x^\gamma$ for all $x>0$ for some $\gamma\in\mb R$. Next, we show that $\gamma = 1$ in the temperate cases $\mathbb{R}_+$ and $\mathbb{R}^{\mathrm{op}}_+$, while $\gamma = 0$ in the tropical cases $\mathbb{TR}_+$ and $\mathbb{TR}^{\mathrm{op}}_+$.
By substituting $x = \frac{1}{2}$, and using the equivalence established in \eqref{eq:x1} along with the (temperate or tropical) additivity, we obtain
\begin{equation}
1 = \Phi(1) =
\begin{cases}
2\,\Phi\!\left(\tfrac{1}{2}\right) = 2^{1-\gamma}, & \quad \mathbb{R}_+,\mathbb{R}_+^{\mathrm{op}}, \\
\Phi\!\left(\tfrac{1}{2}\right) = 2^{-\gamma}, & \quad \mathbb{TR}_+,\mathbb{TR}_+^{\mathrm{op}} \,.
\end{cases}
\end{equation}

This establishes that $\gamma = 1$ in the temperate cases $\mathbb{R}_+$ and $\mathbb{R}^{\mathrm{op}}_+$, and $\gamma = 0$ in the tropical cases $\mathbb{TR}_+$ and $\mathbb{TR}^{\mathrm{op}}_+$. In the above arguments, we implicitly assumed that $P_{Y}(y)>0$. In the situation where $P_Y(y)=0$ we have that $\Phi(P_{X,Y=y}) =0$ since $\Phi(0)=0$. This implies that in the cases $\mathbb{TR}_+$ and $\mathbb{TR}^{\mathrm{op}}_+$, only the $Y$ outcomes such that $P_Y(y)>0$ must be included in the maximization. Substituting the values of $\gamma$ into \eqref{eq:additivity}, we obtain
\begin{equation}\label{eq:general1}
\Phi(P_{XY}) =
\begin{cases}
\sum_{y \in \Y} P_Y(y)\,\Phi\!\left(P_{X|Y=y}\right), & \mathbb{R}_+,\mathbb{R}_+^{\mathrm{op}}, \\
\max \limits_{\substack{y \in \mathcal{Y} \\ P_Y( y) > 0}} \Phi\!\left(P_{X|Y=y}\right), & \mathbb{TR}_+,\mathbb{TR}_+^{\mathrm{op}} \,.
\end{cases}
\end{equation}

 \bigskip
 \bigskip

\noindent \textbf{Step 3: Inner barycenter}. The final step is to determine the form of the homomorphism for a normalized column vector $P_{X|Y=y}$. Using multiplicativity, for two column vectors $P_X$ and $Q_{X'}$, we have that
\begin{equation}
    \Phi(P_X \otimes Q_{X'}) = \Phi(P_X) \Phi(Q_{X'}) \,.
\end{equation}
Furthermore, by monotonicity under mixing channels, we have that for any vector $P_X$ and any doubly stochastic matrix $S$,
\begin{equation}
\begin{cases}
 \Phi(P_X)\leq \Phi(SP_X), & \quad \mathbb{R}_+,\mathbb{TR}_+, \\
\Phi(P_X) \geq \Phi(SP_X), & \quad \mathbb{R}_+^{\mathrm{op}},\mathbb{TR}_+^{\mathrm{op}}\,.
\end{cases}
\end{equation}
Since $\Phi$ has trivial kernel, $\Phi(P_X)>0$ for every probability distribution $P_X$, and hence $P_X\to \log{P_X}$ is well defined.
The latter function is additive under tensor product and monotone under mixing channels. Using Proposition~\ref{proposition:barycenter entropies}, we obtain that the function $\log{\Phi(P_X)}$ is a linear combination of Rényi entropies. This leads to the form
\begin{equation}
\label{eq:FinalForm proof}
\Phi(P_{XY}) =
\begin{cases}
\displaystyle \sum_{y \in \mathcal{Y}} P_Y(y) \exp\Bigg(\int_{[0,+\infty]} H_\alpha(P_{X|Y=y}) \, \mathrm{d}\mu(\alpha) \Bigg), & \mathbb{R}_+,\mathbb{R}_+^{\mathrm{op}}, \\
\displaystyle \max \limits_{\substack{y \in \mathcal{Y} \\ P_Y( y) > 0}} \exp\Bigg(\int_{[0,+\infty]} H_\alpha(P_{X|Y=y}) \, \mathrm{d}\mu(\alpha) \Bigg), & \mathbb{TR}_+,\mathbb{TR}_+^{\mathrm{op}}.
\end{cases}
\end{equation}
where $\mu$ is a positive measure in cases $\mathbb{R}_+$ and $\mathbb{TR}_+$, and it is a negative measure in the opposite cases $\mathbb{R}^{\mathrm{op}}_+$ and $\mathbb{TR}^{\mathrm{op}}_+$.

\bigskip
\bigskip

\noindent \textbf{Step 4: Simplifying $\mathbb{TR}_+$.}
For two probability vectors with the same support, each can be written as a
nontrivial mixture containing the other. Strong quasi-concavity therefore
forces a tropical monotone to take the same value on both vectors. It depends
only on support, so its inner entropy is $H_0$. Consequently,
\begin{equation}
    \Phi(P_{XY}) = \max \limits_{\substack{y \in \mathcal{Y} \\ P_Y( y) > 0}}
    \exp\!\bigl(kH_0(P_{X|Y=y})\bigr), \qquad k>0.
\end{equation}
Propositions~\ref{prop:sufficient-tropical-positive} and
\ref{prop:positive-tropical-necessary} below provide, respectively, the
direct monotonicity proof and the sharp exclusion of every other Rényi
parameter. This completes the proof.
\end{proof}

% \subsection{The degenerate homomorphisms}
% \label{sec: degenerate homomorphisms}
% In this section, we discuss the degenerate monotone homomorphisms of the semiring of conditional majorization.

% \begin{proposition}
% \label{prop: degenerate homomorphisms}
% Let \( S \) denote the semiring of conditional majorization. Then,
% \begin{enumerate}
%     \item In \(\mathbb{R}_+\) and \(\mathbb{R}_+^{\mathrm{op}}\), the only degenerate monotone homomorphism is the total weight \(\|\cdot\|\) introduced in Lemma~\ref{lem: degenerate and zig-zag}.
%     \item In \(\mathbb{TR}_+\), the only degenerate homomorphism is the function that takes the value \(1\) for all nonzero \(P_{XY}\).
% \end{enumerate}
% \end{proposition}

% \begin{proof}
% Let us first start with the cases \(\mathbb{R}_+\) and \(\mathbb{R}_+^{\mathrm{op}}\).
% The monotone homomorphisms $\Phi:\cS\to\mb K\in\{\mb R_+,\mb R_+^{\rm op}\}$ have the form
% \begin{equation}
% \Phi(P_{XY})=\sum_{y\in\Y}P_Y(y)\exp{\left(\int_{[0,+\infty]}H_\alpha(P_{X|Y=y})\d\mu(\alpha)\right)}
% \end{equation}
% where $\mu$ is either a positive or negative finite measure. Hence, the only possibly degenerate homomorphism corresponds to the measure being zero. In this case,
% \begin{equation}
% \Phi(P_{XY})=\sum_{y \in\Y}P_Y(y)=\sum_{x\in\mc X}\sum_{y\in\mc Y}P(x,y)=\|P_{XY}\|\,,
% \end{equation}
% where the total weight of the distribution is represented by the degenerate homomorphism
% $\|\cdot\|$, which characterizes the zig-zag relation $\sim$—was introduced in Lemma~\ref{lem: degenerate and zig-zag}.
% Indeed, the homomorphism is degenerate, since whenever \(P \succeq Q\), it follows that \(\|P\| = \|Q\|\). This is because they are connected by a channel, which necessarily preserves the total weight.


% Let us turn to the cases $\mathbb{TR}_+, \mathbb{TR}_+^{\mathrm{op}} $.
% In this case where $\Phi \colon \cS \to \mathbb{K}$, with $\mathbb{K} \in \{ \mathbb{TR}_+, \mathbb{TR}_+^{\mathrm{op}} \}$, the monotone homomorphisms take the form
% \begin{equation}
% \Phi(P_{XY})=\max_{y \in \Y}\exp{\left(\int_{[0,+\infty]}H_\alpha(P_{X|Y=y})\d\mu(\alpha)\right)} \,,
% \end{equation}
% where $\mu$ is measure. The only degenerate homomorphism corresponds to the measure being zero, i.e.,
% \begin{equation}
% \Phi(P_{XY})=\left\{\begin{array}{ll}
% \max_{y \in \Y}\,1=1,&{\rm when}\ P_{XY}\neq0,\\
% 0,&{\rm when}\ P_{XY}=0.
% \end{array}\right.
% \end{equation}
% This homomorphism is trivial and carries no meaningful structural information; hence, we shall disregard it.

% \end{proof}

% \bigskip

% In the next section, our focus therefore shifts to the study of derivations.

\subsection{The derivations}
\label{sec: derivations}
Besides nondegenerate homomorphisms, the spectral theorem requires monotone
derivations at $\|\cdot\|$. Algebraically these arise from the limiting family
$H_{0,\alpha}$ as $t\to0$; monotonicity restricts the initially available
Rényi orders to the sharp interval $[0,1]$.



\begin{proposition}
\label{prop: monotone derivations}
Let \(S\) be the semiring of conditional majorization. The normalized
extremal monotone derivations at \(\|\cdot\|\), as defined in
Lemma~\ref{lem: degenerate and zig-zag}, are precisely
\begin{equation}
\label{eq:ExtDerivation}
H_{0,\alpha}(P)
= \sum_{y \in \mathcal{Y}} P_Y(y) \, H_\alpha\!\big(P_{X|Y = y}\big)\,,
\end{equation}
for \(\alpha\in[0,1]\). Every normalized monotone derivation is a barycentre
of this family.
\end{proposition}
\begin{proof}
We need to study monotone derivations at $\|\cdot\|$. These are additive maps \(\Delta \colon S \to \mathbb{R}\) that preserve the order, in the sense that \([P] \rgeq [Q]\)—and hence \(P \preceq Q\) for representatives of the corresponding equivalence classes—implies \(\Delta(P) \geq \Delta(Q)\). Moreover, they satisfy the Leibniz rule with respect to the semiring product,
\begin{align}\label{eq:LeibnizRule}
\Delta(P_{XY}\otimes Q_{X'Y'})
= \Delta(P_{XY})\,\|Q_{X'Y'}\|
+ \|P_{XY}\|\,\Delta(Q_{X'Y'}) .
\end{align}

for all $P_{XY},Q_{X'Y'}\in\V$.

The proof is organized into several key steps that lead to the general form of the monotone derivations. Specifically, we proceed as follows:
\begin{enumerate}
    \item Step 1: Additive constant. We begin by showing that the derivations can be taken to satisfy \(\Delta(x) = 0\) for all \(x \in \mathbb{R}_+\).
    \item Step 2: General form. We then exploit the properties of monotone derivations to determine their general form.
    \item Step 3: Extremal derivations. We identify the extremal derivations, which generate the entire set of derivations through convex combinations.
\end{enumerate}


\medskip

\noindent\textbf{Step 1: Additive constant.} We first show that, without loss of generality, the derivation can be assumed to satisfy $\Delta(x) = 0$ for all $x \in \mathbb{R}_+$. To establish this, we prove that from any given derivation $\Delta$, one can construct another derivation $\Delta'$ which satisfies $\Delta'(x) =0$ and hence it is interchangeable with $\Delta$ as defined in Definition 8.2 of \cite{fritz2023abstractII}. Let $\Delta$ be a derivation.  We notice that, for $x,y\in\mb R_+$,
\begin{equation}
\Delta(xy)=\Delta(x)y+x\Delta(y),
\end{equation}
implying that $\delta:=-\Delta|_{\mb R_+}:\mb R_+\to\mb R$ is a derivation.
Define a new map $\Delta' \colon \mathcal{V} \to \mathbb{R}$ by
\begin{equation}
    \Delta'(P) := \Delta(P) - \Delta(\|P\|).
\end{equation}
By construction, $\Delta'(x) = 0$ for all $x \in \mathbb{R}_+$. It therefore suffices to show that $\Delta$ and $\Delta'$ are interchangeable. This will follow if $\Delta'$ is itself a monotone derivation at $\|\cdot\|$. While this is a general fact, we provide a direct proof here for completeness. We have, for all $x,y\in\mb R_+$,
\begin{equation}
\left(\begin{array}{cc}
x&0\\
0&y
\end{array}\right)\succeq\left(\begin{array}{cc}
x+y&0\\
0&0
\end{array}\right)\succeq\left(\begin{array}{cc}
x&0\\
0&y
\end{array}\right),
\end{equation}
by the same argument that we used in the multiplicativity part of Section~\ref{sec: all monotone homomorphisms}. Therefore, we obtain that
\begin{equation}
\Delta(x+y)=\Delta\left[\left(\begin{array}{cc}
x&0\\
0&y
\end{array}\right)\right]=\Delta(x)+\Delta(y)\,.
\end{equation}
In the following, for compactness, we omit the subscripts $X,Y$.
We now have, for all $P,Q\in\V$,
\begin{align}
\Delta'(P\oplus Q)&=\Delta(P\oplus Q)-\Delta(\|P\oplus Q\|)\\
&=\Delta(P)+\Delta(Q)-\Delta(\|P\|+\|Q\|)\\
&=\Delta(P)+\Delta(Q)-\Delta(\|P\|)-\Delta(\|Q\|)\\
&=\Delta'(P)+\Delta'(Q)
\end{align}
and
\begin{align}
\Delta'(P\otimes Q)&=\Delta(P\otimes Q)-\Delta(\|P\otimes Q\|)\\
&=\Delta(P)\|Q\|+\|P\|\Delta(Q)-\Delta(\|P\|\|Q\|)\\
&=\Delta(P)\|Q\|+\|P\|\Delta(Q)-\Delta(\|P\|)\|Q\|-\|P\|\Delta(\|Q\|)\\
&=\Delta'(P)\|Q\|+\|P\|\Delta'(Q),
\end{align}
showing that $\Delta'$ is a derivation at $\|\cdot\|$. To prove monotonicity, we consider $[P]$ and $[Q]$ such that $[P]\rgeq[Q]$. Then, we have
\begin{equation}
\Delta'(P)=\underbrace{\Delta(P)}_{\geq\Delta(Q)}-\Delta(\underbrace{\|P\|}_{=\|Q\|})\geq\Delta(Q)-\Delta(\|Q\|)=\Delta'(Q).
\end{equation}
Thus, $\Delta'$ is also a monotone derivation, and we may replace $\Delta$ with $\Delta'$. This means that we are free to assume that $\Delta(x)=0$ for all $x\in\mb R_+$.

\medskip
\noindent \noindent\textbf{Step 2: General form.}
Next, we derive the algebraic candidate form. On normalized columns,
Proposition~\ref{proposition:barycenter entropies} first gives a finite positive
measure $\tau \colon \mathcal{B}([0,+\infty])\to\R_+$ such that, for a
probability vector $P_X$,
\begin{equation}
\Delta(P_X)=\int_{[0,+\infty]}H_\alpha(P_X)\d\tau(\alpha)\,.
\end{equation}
Combining Lemma~\ref{lemma:shifting} with the additivity and monotonicity of $\Delta$, along with the observations made above, we now obtain the following expression for any joint probability distribution $P_{XY}$
\begin{align}
\Delta(P_{XY})&=\sum_{y\in \Y}\Delta(P_{X,Y=y})=\sum_{y\in \Y}\Delta\big(P_Y(y)P_{X|Y=y}\big)\\
&=\sum_{y\in \Y}\underbrace{\Delta\big(P_Y(y)\big)}_{=0}\|P_{X|Y=y}\|_1+P_Y(y)\Delta(P_{X|Y=y})\\
&=\sum_{y\in \Y}P_Y(y)\Delta(P_{X|Y=y})\\
&=\sum_{y\in \Y}P_Y(y)\int_{[0,+\infty]}H_\alpha(P_{X|Y=y})\d\tau(\alpha).
\end{align}
Hence, we obtain the final form for the derivations
\begin{equation}
\label{eq:derivFinalForm}
    \Delta(P_{XY}) = \sum_{y\in \Y}P_Y(y)\int_{[0,+\infty]}H_\alpha(P_{X|Y=y})\d\tau(\alpha).
\end{equation}
Note that the additivity of $\Delta$ is immediate given the form in \eqref{eq:derivFinalForm}. It is also straightforward to verify that the form above ensures $\Delta$ satisfies the Leibniz rule. Explicitly, if we denote $R_{XX'YY'}:=P_{XY}\otimes Q_{X'Y'}$, the additivity
of the R\'{e}nyi entropies implies that
\begin{align}
\Delta(P_{XY}\otimes Q_{X'Y'})&=\Delta(R_{XX'YY'})\\
&=\sum_{y\in \Y}\sum_{y'\in \Y'}R_{YY'}(y,y')\int_{[0,+\infty]}H_\alpha\big(R_{XX'|YY'=(y,y')}\big)\d\tau(\alpha)\\
&=\sum_{y\in \Y}\sum_{y'\in \Y'}P_Y(y)Q_{Y'}(y')\int_{[0,+\infty]}H_\alpha\big(P_{X|Y=y}\otimes Q_{X'|Y'=y'}\big)\d\tau(\alpha)\\
&=\sum_{y\in Y}\sum_{y'\in \Y'}P_Y(y)Q_{Y'}(y')\int_{[0,+\infty]}H_\alpha\big(P_{X|Y=y}\big)\d\tau(\alpha)\\
&\qquad +\sum_{y\in \Y}\sum_{y'\in \Y'}P_Y(y)Q_{Y'}(y')\int_{[0,+\infty]}H_\alpha\big(Q_{X'|Y'=y'}\big)\d\tau(\alpha)\\
&=\underbrace{\|Q_{Y'}\|_1}_{=\|Q_{X'Y'}\|}\sum_{y\in \Y}P_Y(y)\int_{[0,+\infty]}H_\alpha\big(P_{X|Y=y}\big)\d\tau(\alpha)\\
&\qquad +\underbrace{\|P_Y\|_1}_{=\|P_{XY}\|}\sum_{y'\in \Y'}Q_{Y'}(y')\int_{[0,+\infty]}H_\alpha\big(Q_{X'|Y'=y'}\big)\d\tau(\alpha)\\
&=\Delta(P_{XY})\|Q_{X'Y'}\|+\|P_{XY}\|\Delta(Q_{X'Y'}).
\end{align}
Finally, note that the normalization condition $\Delta\big(P_X\otimes(1/2,1/2)\big)=1$ implies that the positive measure $\tau$ must be a probability measure.


\medskip
\noindent \noindent\textbf{Step 3: Extremal derivations.}
The algebraic extreme candidates are, for $\alpha\in[0,+\infty]$,
\begin{equation}\label{eq:ExtDerivation proof}
H_{0,\alpha}(P)=\sum_{y\in\mc Y} P_Y(y)H_\alpha(P_{X|Y=y})\,.
\end{equation}
From equation~\eqref{eq:derivFinalForm}, we observe that any derivation \(\Delta\) can be written as
\begin{equation}
    \Delta(P) = \int_{[0,+\infty]} H_{0,\alpha}(P)\, d\tau(\alpha)
\end{equation}
for some probability measure \(\tau\). Thus every algebraic candidate is a
barycentre of the $H_{0,\alpha}$. Proposition~\ref{prop:sufficient-derivations}
proves monotonicity for $\alpha\in[0,1]$, while
Proposition~\ref{prop:Derivations} gives an explicit witness for
every $\alpha>1$. The representing measure of a monotone derivation is
therefore supported on $[0,1]$, which proves the statement. Complete
calculations are in Ref.~\cite{rubboli2026parametersFull}; the semiring/channel
lift is checked in Ref.~\cite{rubboli2026parametersLean}.
\end{proof}


% Sections 4 and 5 are maintained directly in this canonical manuscript.
\newcommand{\ParameterMainForms}{the candidate families derived above}
\newcommand{\ParameterMainTemperateForm}{the temperate representation derived above}
\newcommand{\ParameterMainLift}{the convexity-to-monotonicity and semiring-lift results proved in the appendix}
\newcommand{\ParameterMainDerivationLift}{the corresponding derivation results proved in the appendix}
\newcommand{\ParameterMainDerivationRepresentation}{the representation of derivations obtained above}
\newcommand{\ParameterMainSemiring}{the conditional-majorization semiring defined above}
\newcommand{\ParameterMainNormalization}{the normalization fixed above}
\newcommand{\ParameterMainDerivationDefinition}{defined above}
\newcommand{\ParameterMainFinalFamilies}{the temperate, tropical, and derivation candidate families introduced above}
\newcommand{\ParameterSpanOrigin}{is the representation result established above}

\section{Sufficient conditions for the parameters}
\label{sec:sufficient-conditions-measure}

In this section, we establish sufficient conditions for the parameters under which the
homomorphisms and derivations (and, consequently, the associated conditional entropies)
are monotone under conditional majorization.  In the following section, we show that
these conditions are also necessary, without imposing a separate regularity assumption
on the measure.

Throughout Sections~\ref{sec:sufficient-conditions-measure}
and~\ref{sec:necessary-conditions-measures}, every argument headed ``Proof sketch''
is completed analytically in Ref.~\cite{rubboli2026parametersFull}; the corresponding
formal statements and their dependencies are documented in
Ref.~\cite{rubboli2026parametersLean}.  The sketches below record the mathematical
mechanism and the parameter restriction, while those references supply the omitted
estimates and formal verification.

\paragraph{One-column reduction and terminology.}
The semiring, channel preorder, joint-distribution notation, homomorphisms, and
derivations have been defined above.  The remaining parameter question reduces to
the convexity or concavity of a positively homogeneous one-column functional.  We
recall only the endpoint, support, tropical, and discretization terminology used in
the proofs.

The one-column functional for which we need to establish convexity or concavity is
obtained from \ParameterMainTemperateForm:
\begin{align}
\label{eq:TemperateForm}
\Phi(P_{XY})&=\Phi_{t,\tau}(P_{XY})
=\sum_{y\in\Y}P_Y(y)
  \exp\!\left(t\int_{[0,+\infty]}H_\alpha(P_{X|Y=y})\,
  \d\tau(\alpha)\right),
\end{align}
where $\tau$ is a probability measure and
$t\in\R\setminus\{0\}$.  Thus it suffices to study, for
$\vec{x}\in\R^d_{\geq0}\setminus\{0\}$,
\begin{equation}
\label{eq:conditionfunct}
\vec{x}\longmapsto
\|\vec{x}\|_1\exp\!\left(
t\int_{[0,+\infty]}H_\alpha(\hat{x})\,\d\tau(\alpha)
\right),
\qquad
\hat{x}=\frac{\vec{x}}{\|\vec{x}\|_1}.
\end{equation}
Here $d$ is the number of coordinates,
$\|\vec{x}\|_1:=\sum_{i=1}^d x_i$, and $\hat{x}$ is the normalized probability
vector obtained from $\vec{x}$.  For a probability vector $s=(s_i)_i$, the
R\'enyi entropy is
\[
H_\alpha(s):=\frac{1}{1-\alpha}\log\!\left(\sum_i s_i^\alpha\right),
\qquad \alpha\in(0,+\infty)\setminus\{1\},
\]
with endpoint values
\[
H_0(s)=\log|\{i:s_i>0\}|,
\qquad
H_1(s)=-\sum_i s_i\log s_i,
\qquad
H_\infty(s)=-\log\max_i s_i,
\]
where $0\log0:=0$.  We use these endpoint conventions throughout.  The
\emph{support} of a
vector is its set of nonzero coordinates.  For a positive measure $\nu$, its
closed support $\operatorname{supp}(\nu)$ is the set of points every
neighbourhood of which has positive $\nu$-measure.  A point $a$ is an
\emph{atom} when $\nu(\{a\})>0$, and $\delta_a$ is the Dirac probability
measure concentrated at $a$.  A discrete, finitely supported measure is one
concentrated on finitely many parameter values.

The word \emph{temperate} means that $t$ is finite and nonzero.  The positive
and negative \emph{tropical} regimes are the limits $t\to+\infty$ and
$t\to-\infty$, where log-sum-exp becomes a maximum and a minimum,
respectively.  A derivation is the first-order $t\to0$ object.  A
\emph{common discretization} uses the same finitely supported approximation
of the parameter measure at every vector in one convexity or concavity
inequality.  For reference, $F$ is convex when
\[
F(\lambda x+(1-\lambda)z)
\leq \lambda F(x)+(1-\lambda)F(z),
\qquad 0\leq\lambda\leq1,
\]
and concave when this inequality is reversed.  We begin with finite support,
pass to a general measure by a common discretization, and finally take the
tropical and derivation limits.

\subsection{Sufficient conditions for the discrete case}

The discrete case exploits two basic facts.  The first is Minkowski's inequality
\cite[Theorem~9]{bullen2013handbook}.

\begin{lemma}
\label{lem:Minkowski}
Let $\alpha>0$.  The function, defined on
$\R^d_{>0}$,
\begin{equation}
\vec{x}\longmapsto\left(\sum_{i=1}^d x_i^\alpha\right)^{1/\alpha},
\end{equation}
\begin{enumerate}[itemsep=1pt]
\item is convex if $\alpha>1$;
\item is concave if $0<\alpha<1$.
\end{enumerate}
\end{lemma}

The displayed power mean is not defined at $\alpha=0$.  Whenever the entropy
parameter equals $0$, the proof instead uses the endpoint formula for $H_0$.

The second fact is the standard convexity criterion for multivariate monomials
\cite[Lemma~8]{mu21_renyi}; see also Ref.~\cite[Proposition~13]{farooq2024matrix}.

\begin{lemma}
\label{lem:Matrix-majorization}
Let $N\in\mb{N}$ and let
$\vec{\beta}=(\beta_1,\ldots,\beta_N)\in\R^N$ satisfy
$\sum_{\ell=1}^N\beta_\ell=1$.  The function, defined on
$\R^N_{>0}$,
\begin{equation}
\vec{y}\longmapsto\prod_{\ell=1}^N y_\ell^{\beta_\ell},
\end{equation}
\begin{enumerate}[itemsep=1pt]
\item is concave if and only if $\beta_\ell\geq0$ for every $\ell$;
\item is convex if and only if there is a unique index $k$ such that
$\beta_k>0$ and $\beta_\ell\leq0$ for every $\ell\neq k$.
\end{enumerate}
\end{lemma}

Combining these facts gives the finite-support result.

\begin{lemma}
\label{lem:discrete-measure}
Let $\tau$ be a discrete probability measure supported on finitely many points
$\alpha_1<\cdots<\alpha_N$ of $[0,+\infty]$, and let
$t\in\R\setminus\{0\}$.  For every $d\in\mb{N}$, the function
\begin{equation}
\vec{x}\longmapsto\|\vec{x}\|_1
\exp\!\left(t\int_{[0,+\infty]}H_\alpha(\hat{x})\,
\d\tau(\alpha)\right)
\end{equation}
on $\R^d_{\geq0}\setminus\{0\}$ is
\begin{enumerate}
\item concave if $t>0$, $\alpha_1,\ldots,\alpha_N<1$, and
\begin{equation}
\sum_{\ell=1}^N\frac{\alpha_\ell}{1-\alpha_\ell}
\tau(\{\alpha_\ell\})\leq\frac1t;
\end{equation}
\item convex if $t<0$ and either
  \begin{enumerate}[label=(\alph*)]
  \item $\alpha_1,\ldots,\alpha_N\leq1$, or
  \item $\alpha_1,\ldots,\alpha_{N-1}<1$, $\alpha_N>1$, and
  \begin{equation}
  \sum_{\ell=1}^N\frac{\alpha_\ell}{1-\alpha_\ell}
  \tau(\{\alpha_\ell\})\leq\frac1t,
  \end{equation}
  where $\alpha/(1-\alpha)$ has value $-1$ at $\alpha=+\infty$.
  \end{enumerate}
\end{enumerate}
\end{lemma}

\begin{proof}[Proof sketch]
First assume that every atom satisfies
$0<\alpha_\ell<+\infty$ and $\alpha_\ell\neq1$, and set
\begin{equation}
\beta_\ell=t\frac{\alpha_\ell}{1-\alpha_\ell}
\tau(\{\alpha_\ell\}),
\qquad
\bar\beta=1-\sum_{\ell=1}^N\beta_\ell.
\end{equation}
The function factors as
\begin{equation}
\left(\sum_i x_i\right)^{\bar\beta}
\prod_{\ell=1}^N
\left(\sum_i x_i^{\alpha_\ell}\right)^{\beta_\ell/\alpha_\ell}.
\end{equation}
Lemma~\ref{lem:Minkowski} supplies the convexity or concavity of each power mean,
and Lemma~\ref{lem:Matrix-majorization} supplies the joint sign criterion for the
outer monomial.  When $t<0$ and all support lies in $[0,1]$, no moment bound is
needed.  Atoms at $0$, $1$, and $+\infty$ are handled by the endpoint formulas
and endpoint-preserving perturbations or pointwise limits; the branch containing
$\alpha=1$ occurs only in the all-lower negative case.
\end{proof}

\subsection{Sufficient conditions for the temperate homomorphisms - general measure}

We now extend the result to a general probability measure $\tau$.  The signed
integral
\begin{equation}
\label{eq:IntegralCond}
\int_{[0,1)}\frac{\alpha}{1-\alpha}\,\d\tau(\alpha)
+\int_{(1,+\infty]}\frac{\alpha}{1-\alpha}\,\d\tau(\alpha)
\leq\frac1t
\end{equation}
is used only when the two one-sided integrals are finite; at $+\infty$ the
integrand is $-1$.  In particular, Eq.~\eqref{eq:IntegralCond} never hides an
undefined cancellation at $\alpha=1$.
The first integral is the \emph{lower one-sided moment}, and the second is the
\emph{upper one-sided moment}; moment-finiteness means that the relevant
integral is an ordinary finite real number.  We call support in $[0,1]$
\emph{all-lower support}, and support in $(1,+\infty]$ \emph{upper support}.
The exceptional negative branch allows all-lower support together with one
isolated point $\alpha_*>1$.  The \emph{Shannon point}, atom, or branch refers
to $\alpha=1$.

\begin{proposition}
\label{prop:sufficient-conditions}
Let $\tau$ be a probability measure on $[0,+\infty]$ and let
$t\in\R\setminus\{0\}$.  The function in Eq.~\eqref{eq:conditionfunct} is
\begin{enumerate}
\item concave if $t>0$,
$\operatorname{supp}(\tau)\subseteq[0,1]$, $\tau(\{1\})=0$, and
\begin{equation}
\int_{[0,1)}\frac{\alpha}{1-\alpha}\,\d\tau(\alpha)\leq\frac1t;
\end{equation}
\item convex if $t<0$ and either
  \begin{enumerate}[label=(\alph*)]
  \item $\operatorname{supp}(\tau)\subseteq[0,1]$, with no moment-finiteness
  condition, or
  \item there is $\alpha_*\in(1,+\infty]$ such that
  $\operatorname{supp}(\tau)\subseteq[0,1]\cup\{\alpha_*\}$,
  $\tau(\{1\})=0$, the two one-sided integrals in
  Eq.~\eqref{eq:IntegralCond} are finite, and Eq.~\eqref{eq:IntegralCond} holds.
  \end{enumerate}
\end{enumerate}
\end{proposition}

\begin{proof}[Proof sketch]
Use one partition of $[0,+\infty]$ for every vector appearing in the convexity
inequality and push the mass of each cell to a common representative.  On
$[0,1)$ the representatives are chosen from below, so the discrete moment bound
is preserved.  The Shannon branch is kept separate, and in the exceptional
negative case the point $\alpha_*$ is fixed throughout the approximation.  Each
finite approximation satisfies Lemma~\ref{lem:discrete-measure}.

Here a \emph{cell} is one set in the finite partition, and its
\emph{representative} is the single parameter value to which the mass of that
cell is moved.  Calling the discretization common means that these cells and
representatives do not depend on which vector in the inequality is being
evaluated.

For fixed dimension, $0\leq H_\alpha(\hat{x})\leq\log d$, including the endpoint
values.  Dominated convergence therefore gives pointwise convergence of the
entropy integrals and of their exponentials, and a pointwise limit preserves the
required convexity or concavity inequality.  This is the only convergence input
needed here.  Notice especially that the all-lower branch with $t<0$ requires no
integrability of $\alpha/(1-\alpha)$ near $1$.
Here dominated convergence means that one common integrable bound allows the
limit to pass through the integral, and pointwise means at each fixed vector
$\vec{x}$.
\end{proof}

\subsection{Monotonicity under mixing channels}
\label{sec:mixing-channels}

As mentioned at the beginning of this section, monotonicity under mixing channels
acting on the first register $X$ always holds.  Mixing on $X$ increases every
$H_\alpha(P_X)$ and hence the integrated entropy.  Consequently
$\Phi_{t,\tau}$ increases for $t>0$ and decreases for $t<0$; after applying
$(1/t)\log$, the associated conditional entropy increases in both cases.  This
gives the appropriate positive- and negative-homomorphism order directions.
Using \ParameterMainLift, the columnwise concavity or convexity established above
is converted into monotonicity under the conditioning operation and then lifted to
the original conditionally mixing channels.  The derivation case follows from
\ParameterMainDerivationLift.

\subsection{Sufficient conditions for tropical homomorphisms}

We next take the tropical limits.  Monotonicity is preserved under pointwise
limits, so the temperate conditions immediately provide the admissible endpoint
families.

\begin{proposition}
\label{prop:sufficient-tropical}
Let $\tau$ be a probability measure on $[0,+\infty]$.  The function
\begin{equation}
\label{eq:TropicalForm}
P_{XY}\longmapsto
\min_{\substack{y\in\Y\\P_Y(y)>0}}
\int_{[0,+\infty]}H_\alpha(P_{X|Y=y})\,\d\tau(\alpha)
\end{equation}
is monotone under conditionally mixing channels if either
\begin{enumerate}[label=(\alph*)]
\item $\operatorname{supp}(\tau)\subseteq[0,1]$, with no moment-finiteness
condition, or
\item there is $\alpha_*\in(1,+\infty]$ such that
$\operatorname{supp}(\tau)\subseteq[0,1]\cup\{\alpha_*\}$,
$\tau(\{1\})=0$, both one-sided integrals are finite, and
\begin{equation}
\int_{[0,1)}\frac{\alpha}{1-\alpha}\,\d\tau(\alpha)
+\int_{(1,+\infty]}\frac{\alpha}{1-\alpha}\,\d\tau(\alpha)\leq0.
\end{equation}
\end{enumerate}
\end{proposition}

\begin{proof}[Proof sketch]
Use
\begin{align}
\label{eq:limit-tropical}
&\lim_{t\to-\infty}\frac1t\log\!\left(
\sum_{y\in\Y}P_Y(y)
\exp\!\left(t\int H_\alpha(P_{X|Y=y})\,\d\tau(\alpha)\right)
\right)\\
&\hspace{8em}=
\min_{\substack{y\in\Y\\P_Y(y)>0}}
\int H_\alpha(P_{X|Y=y})\,\d\tau(\alpha).
\end{align}
The all-lower case is the limit of Proposition~\ref{prop:sufficient-conditions}(2a).
If the signed moment is strictly negative, the exceptional case holds for all
sufficiently large $|t|$.  At equality, move an arbitrarily small amount of mass
to $\alpha_*$, apply the strict case, and take a final pointwise limit.  This
retains both normalization and the boundary of the parameter range.
\end{proof}

\begin{proposition}
\label{prop:sufficient-tropical-positive}
The function
\begin{equation}
P_{XY}\longmapsto
\max_{\substack{y\in\Y\\P_Y(y)>0}}H_0(P_{X|Y=y})
\end{equation}
is monotone under conditionally mixing channels.
\end{proposition}

\begin{proof}[Proof sketch]
This is the $t\to+\infty$ limit of the positive temperate branch with
$\tau$ concentrated at $\alpha=0$:
\begin{equation}
\lim_{t\to+\infty}\frac1t\log\!\left(
\sum_{y\in\Y}P_Y(y)e^{tH_0(P_{X|Y=y})}\right)
=\max_{\substack{y\in\Y\\P_Y(y)>0}}H_0(P_{X|Y=y}).
\end{equation}
\end{proof}

\begin{remark}
\label{rem:positive-tropical-endpoint}
Unless $\alpha=0$, the positive temperate moment bound is eventually violated as
$t\to+\infty$.  Thus the positive tropical family contains only the
$\alpha=0$ point mass.  In contrast, every point mass is allowed in the negative
tropical family, subject to the two alternatives in
Proposition~\ref{prop:sufficient-tropical}.
\end{remark}

\subsection{Sufficient conditions for derivations}

We now turn to the derivations.

\begin{proposition}
\label{prop:sufficient-derivations}
The function
\begin{equation}
P_{XY}\longmapsto
\sum_{y\in\Y}P_Y(y)H_\alpha(P_{X|Y=y})
\end{equation}
is monotone under conditionally mixing channels if $\alpha\in[0,1]$.
\end{proposition}

\begin{proof}[Proof sketch]
For the point mass at the actual parameter $\alpha$, take $t\to0^-$ in the
all-lower branch of Proposition~\ref{prop:sufficient-conditions}.  The finite
log-sum-exp identity
\begin{equation}
\lim_{t\to0}\frac1t\log\!\left(
\sum_{y\in\Y}P_Y(y)e^{tH_\alpha(P_{X|Y=y})}\right)
=\sum_{y\in\Y}P_Y(y)H_\alpha(P_{X|Y=y})
\end{equation}
passes monotonicity to the limit.  Equivalently, this is the concavity of the
perspective $\vec{x}\mapsto\|\vec{x}\|_1H_\alpha(\hat{x})$ used in the
derivation lift.
The word \emph{perspective} means precisely this positively homogeneous
extension from normalized probability vectors to nonzero nonnegative vectors;
positive homogeneity is the identity $F(s\vec{x})=sF(\vec{x})$ for $s>0$.
\end{proof}

These results identify the admissible homomorphisms and derivations.  With
\ParameterMainNormalization, they give exactly \ParameterMainFinalFamilies.

\section{Necessary conditions for the parameters}
\label{sec:necessary-conditions-measures}

In this section, we show that the sufficient conditions on the parameters derived
in the previous section are also necessary.  No assumption is made that the
measure vanishes at a prescribed rate as $\alpha\to1$, and no global finiteness
condition is placed on the weighted integral.  In the positive branch, concavity
forces the lower one-sided moment to be finite and bounded.  In the exceptional
negative branch, convexity forces both one-sided moments to be finite, whereas
the all-lower negative branch requires no such finiteness.

Our approach consists in identifying explicit perturbation directions and
computing the signs of the second derivatives along them.  Outside the ranges of
Section~\ref{sec:sufficient-conditions-measure}, the resulting curvature has the
wrong sign.  Sharp counterexamples require increasingly high dimensions, but the
underlying idea is the finite-dimensional multivariate monomial
\begin{equation}
\label{eq:multivariate-divergences-function}
 (v_0,\ldots,v_d)\longmapsto
 v_0^{\beta_0}v_1^{\beta_1}\cdots v_d^{\beta_d},
\qquad \beta_0=1-\sum_{i=1}^d\beta_i,
\end{equation}
whose curvature criterion is Lemma~\ref{lem:Matrix-majorization}.  In the
discrete entropy expression, the correspondence is
\begin{equation}
\begin{aligned}
v_0&\longleftrightarrow\sum_jx_j,\\
\beta_i&\longleftrightarrow
\tau(\{\alpha_i\})t\frac{\alpha_i}{1-\alpha_i},\\
v_i&\longleftrightarrow
\left(\sum_jx_j^{\alpha_i}\right)^{1/\alpha_i}.
\end{aligned}
\end{equation}
The latter quantities are coupled through $\vec{x}$.  Blocks whose cardinalities
and weights scale with a dimension parameter make different blocks dominate on
different ranges of $\alpha$.  For example, appending $d$ entries of weight
$1/d$ to a fixed vector $\vec{x}=(x_1,\ldots,x_k)$ gives
\[
\vec{y}=\left(
x_1,\ldots,x_k,
\underbrace{1/d,\ldots,1/d}_{d\text{ entries}}
\right)
\]
and hence
\begin{equation}
\left(\sum_jy_j^\alpha\right)^{1/\alpha}
=\left(\sum_{j=1}^k x_j^\alpha+d^{1-\alpha}\right)^{1/\alpha},
\end{equation}
which diverges for $\alpha<1$ and converges to
$(\sum_{j=1}^k x_j^\alpha)^{1/\alpha}$ for $\alpha>1$.

There is one analytic point that is important even in this abbreviated account.
For each localized path used below, its support is fixed near the perturbation
point and its set of maximizing coordinates remains fixed there.  On these paths the
endpoint-aware kernels and their first two path derivatives extend continuously
through $0$, $1$, and $+\infty$; in particular, the apparent pole at $1$ is
removable.  One differentiates the finite-dimensional kernels before integrating,
then truncates the high-dimensional analysis on the two sides of $1$ and passes
to the endpoint only afterward.  Domination by total variation and monotone
convergence make the required moment bounds conclusions of the argument.  Thus
no derivative is exchanged with an already limiting piecewise envelope, and no
separate finiteness assumption is needed.

An \emph{endpoint-aware kernel} is the function
$\alpha\mapsto H_\alpha(\widehat{\vec{x}+\lambda\vec{v}})$ with the endpoint
values at $0$, $1$, and $+\infty$ retained.  A \emph{removable pole} means
that the apparent denominator $1-\alpha$ is cancelled by a numerator that
also vanishes at $\alpha=1$.  \emph{Endpoint-separated truncation} means
integrating on closed regions below and above $1$ before moving their
boundaries towards $1$.  Domination by total variation means controlling the
absolute value of the kernel by an integrable bound for the positive measure
$|\nu|$ defined next.

For any finite signed measure $\nu$, meaning a finite measure that may take
both positive and negative values, its \emph{total-variation measure}
$|\nu|$ is the positive measure that records the absolute mass of $\nu$, and
we set $\operatorname{supp}(\nu):=\operatorname{supp}(|\nu|)$.  A point $s$
is $|\nu|$-null when $|\nu|(\{s\})=0$; for a positive measure this simply
means $\nu(\{s\})=0$.  Thus a \emph{null point} may still lie in the support:
it has no atom at that single point.  A \emph{null threshold} is a null point
used as the boundary between localization intervals.  A finite measure has at
most countably many atoms, so such thresholds can be chosen between distinct
parameter values.

\paragraph{Explicit witness vectors.}
All counterexamples below are restrictions to affine lines of the form
$\lambda\mapsto\vec{x}+\lambda\vec{v}$.  Here $\lambda\in\R$ is taken small
enough that all coordinates remain positive; $\vec{x}$ is the initial point
at $\lambda=0$, and $\vec{v}$ is its perturbation direction.  The curvature
of a path is its second derivative at $\lambda=0$, while a logarithmic
derivative is a derivative of the logarithm of the positive function being
tested.

Here $d\geq2$ is an auxiliary integer scale parameter tending to infinity.
It is not the actual vector dimension, which is the sum of the block
cardinalities $n_{j,d}$.  A \emph{block} is a group of repeated equal
coordinates, its cardinality is the number of coordinates in that group, and
$\lceil s\rceil$ is the least integer not smaller than $s$.  A block velocity
is the common relative first-order change assigned to one block.  A block
\emph{dominates} at a parameter $\alpha$ when its contribution to
$\sum_i x_i^\alpha$ has the largest power of $d$.  A \emph{localizer} is a
block family arranged so that the dominant block changes with $\alpha$ and
thereby isolates the measure on selected parameter intervals.  A
\emph{witness pair} is the displayed initial point together with its
direction.

We record every coordinate inside parentheses.  For the two-block family, fix
$r>0$ with $r\neq1$.  The number $r$ is the threshold at which the dominant
block changes, and $R:=r+1$ is an auxiliary size exponent.  Put
\begin{equation}
n_{0,d}=\lceil d^R\rceil,
\qquad
n_{1,d}=\lceil d^{R-r}\rceil.
\end{equation}
The initial point and the direction with block velocities $(u_0,u_1)$ are
\begin{align}
\vec{x}_r^{(d)}
&=\left(
\underbrace{1,\ldots,1}_{n_{0,d}},
\underbrace{d,\ldots,d}_{n_{1,d}}
\right),
\label{eq:concise-two-block-point}\\
\vec{v}_r^{(d)}(u_0,u_1)
&=\left(
\underbrace{u_0,\ldots,u_0}_{n_{0,d}},
\underbrace{d u_1,\ldots,d u_1}_{n_{1,d}}
\right).
\label{eq:concise-two-block-direction}
\end{align}
For the three-block family, fix $0<a<b<+\infty$, with $a,b\neq1$.  The two
numbers $a$ and $b$ are the dominant-block crossing thresholds, and
$R:=a+b+1$.  Put
\begin{equation}
n_{0,d}=\lceil d^R\rceil,
\qquad
n_{1,d}=\lceil d^{R-a}\rceil,
\qquad
n_{2,d}=\lceil d^{R-a-b}\rceil.
\end{equation}
Its initial point and direction are
\begin{align}
\vec{x}_{a,b}^{(d)}
&=\left(
\underbrace{1,\ldots,1}_{n_{0,d}},
\underbrace{d,\ldots,d}_{n_{1,d}},
\underbrace{d^2,\ldots,d^2}_{n_{2,d}}
\right),
\label{eq:concise-three-block-point}\\
\vec{v}_{a,b}^{(d)}(u_0,u_1,u_2)
&=\left(
\underbrace{u_0,\ldots,u_0}_{n_{0,d}},
\underbrace{d u_1,\ldots,d u_1}_{n_{1,d}},
\underbrace{d^2u_2,\ldots,d^2u_2}_{n_{2,d}}
\right).
\label{eq:concise-three-block-direction}
\end{align}
Finally, fix a cutoff $c>1$.  It separates the middle parameter range
$1<\alpha<c$ from the upper tail $\alpha>c$ and is not a universal constant.
In each use, $c$ is chosen null for the relevant measure, so no atom
lies on this boundary.  The first two blocks cross at the Shannon point $1$,
and the second and third cross at $c$.  For the rounded Shannon family put
$p=q:=1/2$, $R:=c+2$, and
\begin{equation}
n_{0,d}=\lceil d^R\rceil,
\qquad
n_{1,d}=\lceil d^{R-1}\rceil,
\qquad
n_{2,d}=\lceil d^{R-1-c}\rceil.
\end{equation}
Its initial point and direction are
\begin{align}
\vec{x}_{\mathrm{Sh}}^{(d)}
&=\left(
\underbrace{q,\ldots,q}_{n_{0,d}},
\underbrace{dp,\ldots,dp}_{n_{1,d}},
\underbrace{d^2,\ldots,d^2}_{n_{2,d}}
\right),
\label{eq:concise-Shannon-point}\\
\vec{v}_{\mathrm{Sh}}^{(d)}(h,v)
&=\left(
\underbrace{h/\log d,\ldots,h/\log d}_{n_{0,d}},
\underbrace{-dh/\log d,\ldots,-dh/\log d}_{n_{1,d}},
\underbrace{d^2v,\ldots,d^2v}_{n_{2,d}}
\right).
\label{eq:concise-Shannon-direction}
\end{align}
In the first two families $u_0,u_1,u_2\in\R$ are the block velocities.  In
the Shannon family, $p=q=1/2$ are fixed scalar block weights, while
$h,v\in\R$ control the first-two-block Shannon perturbation and the third
block, respectively; the scalar $v$ is unrelated to a direction vector
$\vec{v}$.  The subscript $\mathrm{Sh}$ stands for Shannon, and ``rounded''
refers to the ceiling signs that turn the desired block sizes into integers.
The velocities will be specified separately for each obstruction.

\subsection{Necessary conditions; temperate case}

The temperate case corresponds to the bulk entropies $H_{t,\tau}$ among
\ParameterMainForms; ``bulk'' here means the finite-temperature, or temperate,
family.  As above, it is sufficient to study
Eq.~\eqref{eq:conditionfunct}.  For notational convenience, combine the weight
$t$ with the probability measure $\tau$ and write
\begin{equation}
\mu:=t\tau.
\end{equation}
A measure $\mu$ is positive when $\mu(E)\geq0$ for every measurable set $E$,
and negative when $\mu(E)\leq0$ for every such $E$.  Since $\mu=t\tau$ is
one-signed, $\mu>0$ means $t>0$, $\mu<0$ means $t<0$, and
$|\mu|=|t|\tau$.  As above, $\operatorname{supp}(\mu)$ means the support of
$|\mu|$.  A lower or upper \emph{truncated coefficient} is a weighted integral
over $[0,a)$ or $(b,+\infty]$, with $a<1<b$; it therefore stays away from the
Shannon point until the two limits are taken separately.
Whenever the two one-sided integrals are finite, we also have the coefficient
\begin{equation}
\label{eq:definition-beta0}
\beta_0=1-
\int_{[0,+\infty]\setminus\{1\}}
\frac{\alpha}{1-\alpha}\,\d\mu(\alpha),
\end{equation}
where the integral is the sum of its parts on $[0,1)$ and $(1,+\infty]$.

Whenever a signed moment is written as a single ordinary real number below, it is
the sum of the two one-sided integrals after their finiteness has been established.
Finiteness is not a blanket hypothesis on $\mu$: before it is known, the argument
works with endpoint-separated truncations.  The localized estimates show that if
convexity leaves any support above $1$, the relevant one-sided moments must be
finite.  If all support lies in $[0,1]$, no such condition is imposed.

We first treat $\mu>0$, and then $\mu<0$.  In the negative case the Shannon atom,
the signed moment, and the number of upper support points are isolated in separate
steps, exactly matching the alternatives in
Proposition~\ref{prop:sufficient-conditions}.

\begin{proposition}
\label{prop:first-measure}
Let $\mu$ be a positive finite measure on $[0,+\infty]$.  The function
\begin{equation}
\vec{x}\longmapsto\|\vec{x}\|_1
\exp\!\left(\int_{[0,+\infty]}H_\alpha(\hat{x})\,
\d\mu(\alpha)\right)
\end{equation}
is concave only if
$\operatorname{supp}(\mu)\subseteq[0,1]$, $\mu(\{1\})=0$, and
\begin{equation}
\int_{[0,1)}\frac{\alpha}{1-\alpha}\,\d\mu(\alpha)\leq1.
\end{equation}
\end{proposition}

\begin{proof}[Proof sketch]
The three exclusions use the finite-dimensional witness families, or
localizers, displayed above.  An atom
$m=\mu(\{1\})>0$ is detected by choosing a $\mu$-null point $c>1$ and using
the complete point-direction pair
\begin{equation}
\left(
\vec{x}_{\mathrm{Sh}}^{(d)},
\vec{v}_{\mathrm{Sh}}^{(d)}(1,0)
\right).
\label{eq:positive-Shannon-witness-pair}
\end{equation}
Its limiting logarithmic derivatives are $m$ and $0$, so the outer exponential
has limiting curvature $m^2>0$.

For the remaining two tests, choose a $\mu$-null threshold $r>0$, put
$R=r+1$, and use the two-block initial point in
Eq.~\eqref{eq:concise-two-block-point}.  If $r>1$ and upper support is present,
the complete pair is
\begin{equation}
\left(
\vec{x}_r^{(d)},
\vec{v}_r^{(d)}(0,1)
\right),
\label{eq:positive-upper-witness-pair}
\end{equation}
and the localized coefficient
\[
A_r=\int_{(r,+\infty]}\frac{\alpha}{1-\alpha}\,\d\mu(\alpha)
\]
is negative.  The limiting first and second logarithmic derivatives are $A_r$
and $-A_r$, hence the curvature is $A_r^2-A_r>0$.  Thus no upper support is
possible.  If $r<1$, use instead the complete pair
\begin{equation}
\left(
\vec{x}_r^{(d)},
\vec{v}_r^{(d)}(1,0)
\right)
\label{eq:positive-lower-witness-pair}
\end{equation}
and write
\[
B_r=\int_{[0,r)}\frac{\alpha}{1-\alpha}\,\d\mu(\alpha).
\]
The limiting derivatives are $B_r$ and $-B_r$, so concavity implies
$B_r(B_r-1)\leq0$ and therefore $B_r\leq1$.
Letting null thresholds increase to $1$ and using monotone convergence proves
the stated moment bound.  Endpoint-separated truncation justifies the limits
without a preliminary finiteness hypothesis.
Monotone convergence here means that the increasing nonnegative truncated
integrands may be passed to their limiting integral.
\end{proof}

Next, let $\mu$ be negative.

\begin{proposition}
\label{prop:no-alpha-one}
Let $\mu$ be a negative finite measure with $\mu(\{1\})\neq0$.  Then the function
\begin{equation}
\vec{x}\longmapsto\|\vec{x}\|_1
\exp\!\left(\int_{[0,+\infty]}H_\alpha(\hat{x})\,
\d\mu(\alpha)\right)
\end{equation}
is convex only if $\operatorname{supp}(\mu)\subseteq[0,1]$.
\end{proposition}

\begin{proof}[Proof sketch]
Assume that both the Shannon atom and upper support are present.  If
$m:=\mu(\{1\})<0$, choose a $|\mu|$-null point $c>1$ such that
$\mu((c,+\infty])<0$, and define the localized upper contribution
\[
A_c:=\int_{(c,+\infty]}
\frac{\alpha}{1-\alpha}\,\d\mu(\alpha)>0.
\]
Use
\begin{equation}
\left(
\vec{x}_{\mathrm{Sh}}^{(d)},
\vec{v}_{\mathrm{Sh}}^{(d)}\!\left(-\frac{A_c}{m},1\right)
\right).
\label{eq:negative-Shannon-witness-pair}
\end{equation}
The first logarithmic derivative then tends to zero and the second tends to
$-A_c<0$.  This is a
genuinely concave direction and contradicts convexity.  Working first on
truncated ranges makes the calculation finite; the endpoint limits are taken only
after the curvature sign is fixed.
\end{proof}

\begin{proposition}
\label{prop:second-measure}
Let $\mu$ be a negative finite measure with $\mu(\{1\})=0$, and suppose that the
function
\begin{equation}
\vec{x}\longmapsto\|\vec{x}\|_1
\exp\!\left(\int_{[0,+\infty]}H_\alpha(\hat{x})\,
\d\mu(\alpha)\right)
\end{equation}
is convex.  Let $0<a<1<b<+\infty$ be $|\mu|$-null thresholds with
$\mu((b,+\infty])<0$, and set
\begin{equation}
B_a=\int_{[0,a)}\frac{\alpha}{1-\alpha}\,\d\mu(\alpha)\leq0,
\qquad
A_b=\int_{(b,+\infty]}\frac{\alpha}{1-\alpha}\,\d\mu(\alpha)>0.
\end{equation}
Thus $B_a$ is the lower truncated coefficient and $A_b$ the upper truncated
coefficient.
Then
\begin{equation}
\label{eq:negative-truncated-moment-necessary}
A_b+B_a\geq1.
\end{equation}
\end{proposition}

\begin{proof}[Proof sketch]
If Eq.~\eqref{eq:negative-truncated-moment-necessary} failed, then
$C:=1-A_b-B_a>0$.  This $C$ is the remaining coefficient of the middle
block, which contains $\alpha=1$.  Use the three-block pair
\begin{equation}
\left(
\vec{x}_{a,b}^{(d)},
\vec{v}_{a,b}^{(d)}\!\left(0,\frac1C,-\frac1{A_b}\right)
\right),
\label{eq:negative-moment-witness-pair}
\end{equation}
where $R=a+b+1$.  The velocity
\[
(u_0,u_1,u_2)=\left(0,\frac1C,-\frac1{A_b}\right)
\]
has limiting first logarithmic derivative $0$ and limiting second derivative
\[
-\frac1C-\frac1{A_b}<0.
\]
This gives a finite-dimensional violation of convexity.  The statement is made
only with truncated coefficients: no singular signed moment is formed before
its two one-sided parts are known to be finite.
\end{proof}

\begin{remark}
\label{rem:normalizing-witnesses}
The points used in these counterexamples need not initially be normalized.  The
function is positively homogeneous on the whole semiring domain, so rescaling by
a positive constant does not change the curvature sign.  Here normalized means
that the coordinates sum to one, and positively homogeneous means
$F(s\vec{x})=sF(\vec{x})$ for $s>0$.  Equivalently, the
counterexample can be embedded into a normalized joint probability matrix,
interpreted in \ParameterMainSemiring, and lifted using \ParameterMainLift.
\end{remark}

\begin{remark}
\label{rem:trace-preserving-witnesses}
A nontraceless direction can likewise be embedded into a trace-preserving
two-column perturbation.  A direction is traceless when $\sum_i v_i=0$;
trace-preserving means that the total mass of all columns remains constant.
Let $\vec{x}'>0$ be a second-column base point and put
$\hat{x}':=\vec{x}'/\|\vec{x}'\|_1$.  If the first column changes as
$\vec{x}+\lambda\vec{v}$, change the second as
\begin{equation}
\vec{x}'-\lambda\hat{x}'\sum_i v_i.
\end{equation}
The opposite sign is essential: it keeps the total weight fixed.  Additivity over
the conditioning register makes the second column affine in $\lambda$, so it
does not alter the curvature obstruction.
\end{remark}

Finally, the upper support cannot contain two distinct points.

\begin{proposition}
\label{prop:third-measure}
Let $\mu$ be a negative finite measure with $\mu(\{1\})=0$, and suppose that the
function
\begin{equation}
\vec{x}\longmapsto\|\vec{x}\|_1
\exp\!\left(\int_{[0,+\infty]}H_\alpha(\hat{x})\,
\d\mu(\alpha)\right)
\end{equation}
is convex.  Then either $\operatorname{supp}(\mu)\subseteq[0,1]$, or there is a
unique $\alpha_*\in(1,+\infty]$ such that
\begin{equation}
\operatorname{supp}(\mu)\subseteq[0,1]\cup\{\alpha_*\}.
\end{equation}
In the second case the two one-sided moments are finite and satisfy
\begin{equation}
\label{eq:negative-moment-necessary}
\int_{[0,1)}\frac{\alpha}{1-\alpha}\,\d\mu(\alpha)
+\int_{(1,+\infty]}\frac{\alpha}{1-\alpha}\,\d\mu(\alpha)\geq1.
\end{equation}
\end{proposition}

\begin{proof}[Proof sketch]
Suppose that $1<\alpha_1<\alpha_2\leq+\infty$ are two support points.  Choose
$|\mu|$-null thresholds
$1<a<\alpha_1<b<\alpha_2$; when $\alpha_2=+\infty$, take any finite null
$b>\alpha_1$ that separates the two support portions.  Define
\[
A_1:=\int_{(a,b)}\frac{\alpha}{1-\alpha}\,\d\mu(\alpha)>0,
\qquad
A_2:=\int_{(b,+\infty]}\frac{\alpha}{1-\alpha}\,\d\mu(\alpha)>0.
\]
These are positive because both $\mu$ and $\alpha/(1-\alpha)$ are negative
above $1$.  The two intervals are localized in different dominant blocks.
The complete witness pair is
\begin{equation}
\left(
\vec{x}_{a,b}^{(d)},
\vec{v}_{a,b}^{(d)}\!\left(0,\frac1{A_1},-\frac1{A_2}\right)
\right).
\label{eq:negative-two-upper-witness-pair}
\end{equation}
Its direction cancels the limiting first derivative, while the
limiting second derivative is
\begin{equation}
-\frac1{A_1}-\frac1{A_2}<0.
\end{equation}
This contradicts convexity.  Shrinking the two localization windows shows that
the upper support has at most one point.  The finiteness and signed inequality in
the exceptional branch now follow from Proposition~\ref{prop:second-measure}:
the single upper atom gives a finite upper moment, its truncated inequality
bounds the lower terms, and monotone convergence gives lower-moment finiteness.
Only then are the two one-sided integrals combined.
\end{proof}

Thus the last two propositions cover all negative measures without a global
finiteness assumption; the all-lower branch still carries no moment condition.

\subsection{Necessary conditions for tropical homomorphisms}
\label{sec:necessary-tropical}

The tropical counterexamples use the same finite witness families as the
temperate proof, but quasiconvexity permits a second-derivative test only after
the first derivative has been made exactly zero at the finite dimension under
consideration.
Here quasiconvexity means
\[
G(\lambda x+(1-\lambda)z)
\leq\max\{G(x),G(z)\},
\qquad 0\leq\lambda\leq1.
\]
A direction is \emph{stationary} when its first derivative is zero.  A
\emph{transverse direction} is a correction direction whose first derivative
is nonzero, so a small multiple of it can impose exact stationarity.

\begin{proposition}
\label{prop:negative-tropical-necessary}
Let $\tau$ be a probability measure on $[0,+\infty]$.  If the function in
Eq.~\eqref{eq:TropicalForm} is monotone under conditionally mixing channels,
then either
\begin{enumerate}[label=(\alph*)]
\item $\operatorname{supp}(\tau)\subseteq[0,1]$; or
\item there is a unique $\alpha_*\in(1,+\infty]$ such that
\[
\operatorname{supp}(\tau)\subseteq[0,1]\cup\{\alpha_*\},
\qquad \tau(\{1\})=0,
\]
both one-sided integrals
\[
\int_{[0,1)}\frac{\alpha}{1-\alpha}\,\d\tau(\alpha),
\qquad
\int_{(1,+\infty]}\frac{-\alpha}{1-\alpha}\,\d\tau(\alpha)
\]
are finite, and
\begin{equation}
\label{eq:negative-tropical-necessary-moment}
\int_{[0,+\infty]\setminus\{1\}}
\frac{\alpha}{1-\alpha}\,\d\tau(\alpha)\leq0.
\end{equation}
\end{enumerate}
\end{proposition}

\begin{proof}[Proof sketch]
A Shannon localizer first excludes $\tau(\{1\})>0$ whenever upper support is
present.  Set $\mu:=-\tau$.  Then the signed Shannon atom is
$m:=\mu(\{1\})=-\tau(\{1\})<0$.  Choose a $|\mu|$-null point $c>1$ with
$\mu((c,+\infty])<0$ and define
\[
A_c:=\int_{(c,+\infty]}
\frac{\alpha}{1-\alpha}\,\d\mu(\alpha)>0.
\]
Use the initial point and the two displayed directions
\begin{equation}
\left(
\vec{x}_{\mathrm{Sh}}^{(d)},
\vec{v}_{\mathrm{Sh}}^{(d)}\!\left(-\frac{A_c}{m},1\right)
\right),
\qquad
\left(
\vec{x}_{\mathrm{Sh}}^{(d)},
\vec{v}_{\mathrm{Sh}}^{(d)}(0,1)
\right),
\label{eq:tropical-Shannon-point-directions}
\end{equation}
corresponding to the block-velocity vectors
$u=(-A_c/m,1)$ and $e=(0,1)$.  The limiting first derivative along $u$ is
zero, the second derivative is $-A_c<0$, and the limiting first derivative
along $e$ is $A_c\neq0$.

If two separated upper-support portions remain, choose support points
$1<\alpha_1<\alpha_2\leq+\infty$ and $|\mu|$-null thresholds
$1<a<\alpha_1<b<\alpha_2$, with the same finite-$b$ interpretation as above
when $\alpha_2=+\infty$.  Define
\[
A_1:=\int_{(a,b)}\frac{\alpha}{1-\alpha}\,\d\mu(\alpha)>0,
\qquad
A_2:=\int_{(b,+\infty]}\frac{\alpha}{1-\alpha}\,\d\mu(\alpha)>0.
\]
These coefficients give the initial point and directions
\begin{equation}
\left(
\vec{x}_{a,b}^{(d)},
\vec{v}_{a,b}^{(d)}\!\left(0,\frac1{A_1},-\frac1{A_2}\right)
\right),
\qquad
\left(
\vec{x}_{a,b}^{(d)},
\vec{v}_{a,b}^{(d)}(0,1,0)
\right),
\label{eq:tropical-two-upper-point-directions}
\end{equation}
corresponding to
$u=(0,1/A_1,-1/A_2)$ and $e=(0,1,0)$.  The limiting first derivative along
$u$ is zero, its second derivative is $-1/A_1-1/A_2<0$, and the first
derivative along $e$ tends to $A_1\neq0$.  If
Eq.~\eqref{eq:negative-tropical-necessary-moment} fails in the exceptional
branch, choose $|\mu|$-null thresholds $0<a<1<b<\alpha_*$ close enough to
$1$ that
\[
B_a:=\int_{[0,a)}\frac{\alpha}{1-\alpha}\,\d\mu(\alpha)\leq0,
\qquad
A_b:=\int_{(b,+\infty]}\frac{\alpha}{1-\alpha}\,\d\mu(\alpha)>0
\]
satisfy $A_b+B_a<0$, and define
\[
C:=-A_b-B_a>0.
\]
This tropical coefficient has no added $1$: unlike the temperate function,
the tropical one-column function has no outer $\|\vec{x}\|_1$ term.  The
coefficients $A_b,C>0$ give
\begin{equation}
\left(
\vec{x}_{a,b}^{(d)},
\vec{v}_{a,b}^{(d)}\!\left(0,\frac1C,-\frac1{A_b}\right)
\right),
\qquad
\left(
\vec{x}_{a,b}^{(d)},
\vec{v}_{a,b}^{(d)}(0,0,1)
\right),
\label{eq:tropical-moment-point-directions}
\end{equation}
corresponding to $u=(0,1/C,-1/A_b)$ and $e=(0,0,1)$.  Here the limiting first
derivative along $u$ is zero, its second derivative is $-1/C-1/A_b<0$, and the
first derivative along $e$ tends to $A_b\neq0$.

For each of the three witnesses, let $u$ and $e$ have the values displayed in
the relevant formula.  At every sufficiently large finite dimension, replace
$u$ by $u+\theta_de$, where $\theta_d\in\R$ is the scalar correction chosen
so that the first derivative is exactly zero.  The transverse first derivative
tends to a nonzero number,
whereas the first derivative along $u$ tends to zero, so $\theta_d\to0$ and the
corrected second derivative remains strictly negative.  Let $g_d(\lambda)$ be
the negative tropical one-column function restricted to this corrected affine
line.  Then
\begin{equation}
g_d'(0)=0,
\qquad
g_d''(0)<0.
\label{eq:tropical-exact-negative-curvature}
\end{equation}
Taylor expansion gives, for every sufficiently small $\varepsilon>0$,
\begin{equation}
g_d(-\varepsilon)<g_d(0),
\qquad
g_d(\varepsilon)<g_d(0).
\label{eq:tropical-endpoint-strict-inequalities}
\end{equation}
Quasiconvexity at the midpoint would instead require
\begin{equation}
g_d(0)\leq
\max\{g_d(-\varepsilon),g_d(\varepsilon)\},
\label{eq:tropical-quasiconvex-midpoint}
\end{equation}
which is impossible.  Thus all three negative second derivatives give genuine
finite-dimensional tropical counterexamples.  Support isolation makes the
upper moment finite; monotone passage through lower truncations then proves
lower-moment finiteness.
\end{proof}

\begin{proposition}
\label{prop:positive-tropical-necessary}
Let $\tau$ be a probability measure on $[0,+\infty]$.  If
\begin{equation}
\label{eq:general-positive-tropical-candidate}
P_{XY}\longmapsto
\max_{\substack{y\in\Y\\P_Y(y)>0}}
\int_{[0,+\infty]}H_\alpha(P_{X|Y=y})\,\d\tau(\alpha)
\end{equation}
is monotone under conditionally mixing channels, then $\tau=\delta_0$.
Consequently, Eq.~\eqref{eq:general-positive-tropical-candidate} is precisely
the function in Proposition~\ref{prop:sufficient-tropical-positive}.
\end{proposition}

\begin{proof}[Proof sketch]
For a finite probability vector $s$, write
\begin{equation}
A_\tau(s)=\int_{[0,+\infty]}H_\alpha(s)\,\d\tau(\alpha).
\label{eq:concise-positive-tropical-integrated-entropy}
\end{equation}
For probability vectors $p,q$ of the same dimension, positive-tropical
monotonicity gives the max-quasiconcavity inequality
\begin{equation}
A_\tau(\lambda p+(1-\lambda)q)
\geq\max\{A_\tau(p),A_\tau(q)\},
\qquad 0<\lambda<1.
\label{eq:concise-strong-max-quasiconcavity}
\end{equation}
A fixed-support simplex face is the set of probability vectors with one fixed
pattern of allowed nonzero coordinates; its relative interior consists of
the vectors strictly positive on every allowed coordinate.  Applying the
displayed inequality in both directions on such a face makes $A_\tau$ constant
on its relative interior.  Apply it to the binary points
\[
p=\left(\frac13,\frac23\right),
\qquad
u=\left(\frac12,\frac12\right).
\]
Here $p$ is a probability vector, unrelated to the scalar Shannon block weight
$p=1/2$ used above, and $u$ is the uniform binary vector.
They have the same support, so $A_\tau(p)=A_\tau(u)=\log2$.  Moreover,
$H_0(p)=\log2$ and $H_\alpha(p)<\log2=H_\alpha(u)$ for every
$\alpha\in(0,+\infty]$.  Hence
\begin{equation}
0=\int_{[0,+\infty]}
\bigl(\log2-H_\alpha(p)\bigr)\,\d\tau(\alpha).
\label{eq:concise-positive-tropical-zero-gap}
\end{equation}
The integrand is nonnegative everywhere and strictly positive on
$(0,+\infty]$, so $\tau((0,+\infty])=0$ and $\tau=\delta_0$.
\end{proof}

\subsection{Necessary conditions for derivations}
\label{sec:necessary-derivations}

Next, we turn to the necessary conditions for the extremal derivations.  Under
\ParameterMainDerivationRepresentation, these are
\begin{equation}
H_{0,\alpha}(P_{XY})=
\sum_{y\in\Y}P_Y(y)H_\alpha(P_{X|Y=y}).
\end{equation}
It is enough to work directly at the measure level.  Thus consider
\begin{equation}
\vec{x}\longmapsto\|\vec{x}\|_1
\int_{[0,+\infty]}H_\alpha(\hat{x})\,\d\tau(\alpha),
\end{equation}
for a finite positive measure $\tau$.  Point masses recover the extremal
derivations used in \ParameterMainDerivationLift.
Here \emph{extremal} means the point-mass case $\tau=\delta_\alpha$; a general
represented derivation is a positive integral mixture of these cases.

\begin{proposition}
\label{prop:measure-derivation}
Let $\tau$ be a finite positive measure.  The function
\[
\vec{x}\longmapsto\|\vec{x}\|_1
\int_{[0,+\infty]}H_\alpha(\hat{x})\,\d\tau(\alpha)
\]
is concave only if
$\operatorname{supp}(\tau)\subseteq[0,1]$.  In particular,
$\vec{x}\mapsto\|\vec{x}\|_1H_\alpha(\hat{x})$ is concave only if
$\alpha\in[0,1]$.
\end{proposition}

\begin{proof}[Proof sketch]
If the support meets $(1,+\infty]$, choose a $\tau$-null threshold $r>1$ with
$\tau((r,+\infty])>0$, put $R=r+1$, and define afresh
\[
A_r:=\int_{(r,+\infty]}
\frac{\alpha}{1-\alpha}\,\d\tau(\alpha)<0.
\]
Use the complete pair
\begin{equation}
\left(
\vec{x}_r^{(d)},
\vec{v}_r^{(d)}(0,1)
\right).
\label{eq:derivation-upper-witness-pair}
\end{equation}
Its localized upper
coefficient $A_r$ is negative, while the second derivative of the normalized
derivation converges to $-A_r>0$.  This is a finite-dimensional convex direction
and contradicts concavity, including when the upper support point is
$+\infty$.
\end{proof}

Combining \ParameterMainDerivationRepresentation\ with
Proposition~\ref{prop:sufficient-derivations} and the semiring lift, gives the
final statement.

For clarity, $S$ is \ParameterMainSemiring; its elements are represented by
unnormalized joint probability matrices, and its addition
and multiplication are direct sum and tensor product.  The distinguished norm
$\|P\|$ is the total weight, namely the sum of all entries of $P$.  A
derivation $\Delta$ at this norm is additive and satisfies
\[
\Delta(P\otimes Q)
=\Delta(P)\|Q\|+\|P\|\Delta(Q).
\]

\begin{proposition}
\label{prop:Derivations}
The set of derivations of the conditional majorization semiring $S$ at
$\|\cdot\|$ is spanned by the quantities $H_{0,\alpha}$
\ParameterMainDerivationDefinition, with $\alpha\in[0,1]$.
\end{proposition}

The span assertion in Proposition~\ref{prop:Derivations} \ParameterSpanOrigin.
Here ``spanned'' means generated by positive
barycentres: a represented derivation has the form
$\int H_{0,\alpha}\,\d\tau(\alpha)$ for a finite positive measure $\tau$;
under \ParameterMainNormalization, $\tau$ is a probability measure.
The sufficiency and curvature arguments in this section supply the exact
restriction $\alpha\in[0,1]$ for its extremal candidates.

Together, the temperate, tropical, and derivation statements give the exact
parameter ranges for \ParameterMainFinalFamilies.  Together with the preceding
representation results, this completes the parameter classification.




\section{Large sample and catalytic conditional majorization}
\label{sec: large sample}
In this section, we present the large-sample and catalytic results that provide
the ordering input for the general form of a conditional entropy in
Theorem~\ref{th:barycenter conditional entropies}.
In particular, these results follow directly from Theorem~\ref{thm:Fritz2022} on large-sample and catalytic transformations for general semirings, applied to the semiring of conditional majorization, together with the characterization of monotone homomorphisms and derivations established in Section~\ref{sec: semiring of conditional entropies}.
As shown in the following section (see Theorem~\ref{th:barycenter conditional entropies}), these homomorphisms and derivations coincide, up to a constant and a logarithmic factor, with the extremal conditional entropies that generate the entire family of conditional entropies through convex combinations.
Since conditional entropies, according to Definition~\ref{def: conditional entropy}, require proper normalization and additivity under tensor products, we reformulate the monotone homomorphisms from Proposition~\ref{prop: monotone homomorphisms} and the derivations from Proposition~\ref{prop: monotone derivations} as follows.



\begin{definition}
\label{def: extremal entropies}
Let $P_{XY}$ be a probability distribution, $t\in\R\setminus\{0\}$, and $\tau:\mc B\big([0,+\infty]\big)\to\mb R$ be a probability measure. Then, we define
\begin{align}
    &H_{t,\tau}(X|Y)_P = \frac{1}{t}\log \left(\sum_{y \in \Y} P_Y(y)\exp{\left(t\int_{[0,+\infty]}H_\alpha(P_{X|Y=y})\d\tau(\alpha)\right)}\right) \,, \label{eq:TempEnt}\\
    & H_{-\infty,\tau}(X|Y)_P =  \min \limits_{\substack{y \in \mathcal{Y} \\ P_Y( y) > 0}}\int_{[0,+\infty]}H_\alpha(P_{X|Y=y})\d\tau(\alpha) \,, \label{eq:TropEnt}\\
    & H_{0,\alpha}(X|Y)_P =  \sum_{y \in \Y} P_Y(y)H_\alpha(P_{X|Y=y})\,,\label{eq:DerivEnt} \\
    & H_{+\infty,0}(X|Y)_P =  \max \limits_{\substack{y \in \mathcal{Y} \\ P_Y( y) > 0}} H_0(P_{X|Y=y}) \,. \label{eq:TropEntno}
\end{align}
\end{definition}
Throughout this manuscript, we refer to \(H_{t,\tau}\) as the bulk conditional entropies, as they are defined for finite values of the parameter \(t\). This distinguishes them from the entropies \(H_{-\infty,\tau}\), \(H_{+\infty,0}\), and \(H_{0,\alpha}\), which arise as limiting cases of $H_{t,\tau}$ corresponding to \(t \to -\infty\), \(t \to +\infty\), and \(t \to 0\), respectively.



As we prove later in Theorem~\ref{th:barycenter conditional entropies}, the above quantities encompass all extremal conditional entropies for appropriate values of $t$ and $\tau$. Not every parameter choice is admissible: for some choices the resulting quantity violates the axioms of Definition~\ref{def: conditional entropy}, in particular monotonicity under conditionally mixing channels. The exact admissible regimes are established in Sections~4 and~5, proved in full analytic detail in Ref.~\cite{rubboli2026parametersFull}, and formally verified in Ref.~\cite{rubboli2026parametersLean}. Consequently, we introduce the following sets of parameters and measures.



\begin{definition}
\label{def: D}
Let \(\D_{\rm{bulk}}\) be the set of pairs $(t,\tau)$ consisting of $t\in\R\setminus\{0\}$ and probability measures $\tau:\mc B\big([0,+\infty]\big)\to\mb R$ satisfying one of the following conditions
\begin{enumerate}[itemsep=0pt,topsep=5pt]
    \item $t > 0$, $\;\textup{supp}(\tau) \subseteq [0,1]$, $\tau(\{1\})=0$ and
    \begin{equation}\label{eq:Itaas}
    \int_{[0,+\infty]\setminus\{1\}}\frac{\alpha}{1-\alpha}\,\mathrm{d}\tau(\alpha)\leq\frac{1}{t}
    \end{equation}
    or
    \item $t < 0$ and
    \begin{itemize}
    \item[(a)] $\;\textup{supp}(\tau) \subseteq [0,1]$ or
    \item[(b)] $\;\operatorname{supp}(\tau) \subseteq [0,1] \cup \{\alpha_*\}$, $\tau(\{1\})=0$ for some $\alpha_* \in (1, +\infty]$, both one-sided integrals
    \begin{equation*}
    \int_{[0,1)}\frac{\alpha}{1-\alpha}\,\mathrm{d}\tau(\alpha)
    \quad\text{and}\quad
    \int_{(1,+\infty]}\frac{-\alpha}{1-\alpha}\,\mathrm{d}\tau(\alpha)
    \end{equation*}
    are finite, and \eqref{eq:Itaas} is satisfied.
    \end{itemize}
\end{enumerate}
\end{definition}

\begin{definition}
\label{def: D_infty}
 Let \(\D_{-\infty}\) be the set of probability measures $\tau:\mc B\big([0,+\infty]\big)\to\mb R_+$ satisfying one of the following conditions
    \begin{enumerate}[itemsep=0pt,topsep=5pt]
    \item  $\;\textup{supp}(\tau) \subseteq [0,1]$.
    \item $\;\operatorname{supp}(\tau) \subseteq [0,1] \cup \{\alpha_*\}$, $\tau(\{1\})=0$ for some $\alpha_* \in (1, +\infty]$, both one-sided integrals
    \begin{equation*}
    \int_{[0,1)}\frac{\alpha}{1-\alpha}\,\mathrm{d}\tau(\alpha)
    \quad\text{and}\quad
    \int_{(1,+\infty]}\frac{-\alpha}{1-\alpha}\,\mathrm{d}\tau(\alpha)
    \end{equation*}
    are finite, and
    \begin{align}
        \;\displaystyle\int_{[0,+\infty]\setminus\{1\}} \frac{\alpha}{1 - \alpha} \, \mathrm{d}\tau(\alpha) \leq 0\,.
    \end{align}
\end{enumerate}
\end{definition}
We note that, whenever it is imposed, the condition $\tau(\{1\})=0$
excludes an atom at $\alpha=1$, while still allowing the measure to
accumulate at $1$. This differs from requiring
$\text{supp}(\tau)\subseteq[0,1)$, which would also exclude such accumulation,
since $1$ would then belong to the topological support.


Next, we present the large-sample and catalytic results for transforming joint probability distributions under conditionally mixing operations, relabeling, and embedding, which jointly define the conditional majorization order.

\begin{theorem}
\label{th: large sample and catalytic}
    Let $P_{XY}$ and $Q_{X'Y'}$ be joint probability distributions. If
    \begin{align}
    & H_{t,\tau}(X|Y)_P < H_{t,\tau}(X'|Y')_Q, \qquad\qquad\qquad   \forall (t,\tau) \in \D_{\rm{bulk}},\label{eq:TempCond} \\
    &H_{-\infty,\tau}(X|Y)_P < H_{-\infty,\tau}(X'|Y')_Q, \; \quad\qquad \;\;\, \forall \tau\in \D_{-\infty},\label{eq:TropCond} \\
    &H_{0,\alpha}(X|Y)_P<H_{0,\alpha}(X'|Y')_Q,  \quad\,\quad\qquad\quad\;\; \forall \alpha\in[0,1]\,,\label{eq:DerivCond} \\
    &H_{+\infty,0}(X|Y)_P < H_{+\infty,0}(X'|Y')_Q\,,
    \end{align}
    then
   \begin{enumerate}[itemsep=0pt]
    \item for sufficiently large \(n \in \mathbb{N}\), it holds that $P_{XY}^{\otimes n} \succeq Q_{X'Y'}^{\otimes n}$.
    \item there exists a joint probability distribution \(R_{X''Y''}\) such that
    \begin{equation}
        P_{XY} \otimes R_{X''Y''} \succeq Q_{X'Y'} \otimes R_{X''Y''}.
    \end{equation}
\end{enumerate}
\end{theorem}
\begin{proof}
The result follows directly from an application of the general large-sample theorem for preordered semirings stated in Theorem~\ref{thm:Fritz2022} along with the specific characterization of monotone homomorphisms and derivations for the conditional majorization semiring in Proposition~\ref{prop: monotone homomorphisms} and~\ref{prop: monotone derivations}. Indeed, the entropies appearing in the assumptions coincide with the homomorphisms and derivations up to additive constants and logarithmic factors, which preserve the monotonicity properties. Especially for the derivations, note that, according to Proposition \ref{prop: monotone derivations}, all derivations at $\|\cdot\|$ are spanned by $H_{0,\alpha}$, $\alpha\in[0,1]$, so having strict inequalities over all these derivations is equivalent to having the strict inequalities of \eqref{eq:DerivCond}. Hence, by appropriately accounting for the direction of the inequalities, the assumptions of the theorems coincide with those required in Theorem~\ref{thm:Fritz2022} on the ordering of homomorphisms and derivations.

In addition, according to Theorem~\ref{thm:Fritz2022}, we need to verify that the input distribution $Q_{X'Y'}$ is power universal. We note that here that \(Q_{X'Y'}\) is the input distribution, rather than \(P_{XY}\), since the preorder of the semiring \(\rgeq\) is reversed relative to the conditional majorization \(\succeq\) according to the definition of the preorder of the conditional majorization semiring in item~\ref{item: preorder} of Section~\ref{sec: semiring of conditional entropies}.
The fact that $Q_{X'Y'}$ is power universal is already guaranteed by the above assumptions on the strict ordering of entropies. Indeed, if $Q_{X'Y'}$ were not power universal, according to Proposition~\ref{prop: power universals}, there would exist a column of $Q_{X'Y'}$ with exactly one nonzero entry.
In this case, the Hayashi conditional entropy, which is an element of the set \(\mathcal{D}_{-\infty}\)—a well-known fact in the literature (see, e.g., \cite{gour2024inevitability})— would be zero for $Q_{X'Y'}$. Indeed,
\begin{equation}
    H^{\rm{H}}_{\infty}(X'|Y')_Q = \min_{y \in \mathcal{Y}'}  H_{\infty}\bigl(Q_{X'|Y'=y'}\bigr) \, = 0.
\end{equation}
Since $ H^{\rm{H}}_{\infty}(X|Y)_P\geq0$, the strict inequalities assumed in the theorem imply that $Q_{X'Y'}$ must be power universal.
\end{proof}

 \section{The set of conditional entropies}
\label{sec: set of conditional entropies}
In this section, we present the main result of our work, namely the general characterization of conditional entropies in Theorem \ref{th:barycenter conditional entropies}. Specifically, we show that any conditional entropy defined as a function satisfying the axioms in Definition~\ref{def: conditional entropy} can be expressed as a (generalized) convex combination of the entropies $H_{t,\tau}$ and its limiting forms $H_{0,\alpha}$, $H_{-\infty,\tau}$ and $H_{+\infty,0}$ defined in~\ref{def: extremal entropies}. In this sense, the conditional entropies $H_{t,\tau},H_{0,\alpha},H_{-\infty,\tau}$ and $H_{+\infty,0}$ constitute the extremal points of the set of the conditional entropies; for the extremality, see Theorem 14 of \cite{haapasalo_inprep}. In Section~\ref{sec: literature}, we review more in depth how the conditional entropies discussed in the literature and also mentioned in the introduction relate to our general framework.


%While one could formulate the result more broadly by allowing convex combinations over a larger family—namely, by asserting that any conditional entropy can be expressed as a convex combination of extremal conditional entropies for some choice of parameters—this description is redundant. It therefore suffices to restrict attention to the entropies $H_{t,\tau}$, $H_{0,\alpha}$, and $H_{-\infty,\tau}$ with parameters in the sets $\mathcal{D}_{\mathrm{bulk}}$, $\mathcal{D}_{0}$, and $\mathcal{D}_{-\infty}$, respectively, which precisely characterize the parameter regimes for which the corresponding extremal conditional entropies are monotone under conditionally mixing channels (see Definition~\ref{def: true sets}). Since the remaining axioms are readily verified, these entropies therefore satisfy all axioms required of a conditional entropy (see Definition~\ref{def: conditional entropy}).




%Concerning the proof strategy, the result follows directly from Theorem~\ref{th: large sample and catalytic}, which characterizes the monotones providing sufficient conditions for transformations between multiple copies of joint probability distributions under conditionally mixing channels (large-sample conditional majorization). The argument is then completed by invoking standard results that connect large-sample characterizations with barycentric decompositions. Specifically, these results imply essentially that any quantity that is both monotone and additive can be written as a convex combination of the functions that provide sufficient conditions for the corresponding large-sample transformations (see, e.g.,~\cite{mu21_renyi, haapasalo_inprep}).

\subsection{The general form of conditional entropies}

Let us note that the set of conditional entropies characterized in Definition \ref{def: conditional entropy} is convex. This means that, whenever $\mb H_1$ and $\mb H_2$ are conditional entropies (i.e., satisfy the postulates of Definition \ref{def: conditional entropy}) and $\lambda\in[0,1]$, then also $\mb H=\lambda\mb H_1+(1-\lambda)\mb H_2$,
\begin{equation}
\mb H(X|Y)_P=\lambda\mb H_1(X|Y)_P+(1-\lambda)\mb H_2(X|Y)_P\qquad\forall\, P=P_{XY}\,,
\end{equation}
is a conditional entropy. Especially all convex combinations of the conditional entropies $H_{t,\tau}$, $(t,\tau)\in\mc D_{\rm bulk}$, and $H_{-\infty,\tau}$, $\tau\in\mc D_{-\infty}$, and their pointwise limits are also conditional entropies. This fact also extends to the generalized convex combinations of these quantities, i.e.,\ barycentres over the set of these specific conditional entropies with respect to any (inner and outer regular) probability measure. For the latter statement, we need the fact that the set of these (extremal) conditional entropies is a compact Hausdorff space with respect to the topology of pointwise convergence \cite[Proposition 8.5]{fritz2023abstractII}. The following theorem characterizing the total set of conditional entropies as described in Definition \ref{def: conditional entropy} tells us that any conditional entropy can be expressed as a generalized convex combination of the special conditional entropies characterized in this work.


We recall that $H_{t,\tau}$, $H_{-\infty,\tau}$, $H_{0,\alpha}$, and $H_{+\infty,0}$ are introduced in Definition~\ref{def: extremal entropies}, while the sets $\mathcal{D}_{\mathrm{bulk}}$ and $\mathcal{D}_{-\infty}$ are defined in Definitions~\ref{def: D} and~\ref{def: D_infty}, respectively.
\begin{theorem}[General form of conditional entropies]
\label{th:barycenter conditional entropies}
Let $\H$ be a conditional entropy in the sense of Definition~\ref{def: conditional entropy}.
Then $\H$ can be represented as a generalized convex combination of the conditional entropies $H_{t,\tau}, H_{-\infty,\tau}, H_{0,\alpha}$ and $H_{+\infty,0}$, with parameters $(t,\tau)\in\mathcal{D}_{\mathrm{bulk}}$, $\tau\in\mathcal{D}_{-\infty}$,  and $\alpha\in[0,1]$, respectively, for the first three families.
\end{theorem}
\begin{proof}
We do not provide all details of the proof here, as this result is proven analogously to Proposition \ref{proposition:barycenter entropies} and it is a special case of Theorem 7 of \cite{haapasalo_inprep}. However, we give an overview of the main steps for completeness. Let us denote by $\mc H$ the set of conditional entropies $H_{t,\tau}$, $H_{-\infty,\tau}$, $H_{0,\alpha}$, and $H_{+\infty,0}$ as defined in \eqref{eq:TempEnt}, \eqref{eq:TropEnt}, and \eqref{eq:DerivEnt} of Definition~\ref{def: extremal entropies}, where the parameters and measures satisfy the required monotonicity conditions. According to Proposition 8.5 of \cite{fritz2023abstractII}, this set is a compact Hausdorff space when equipped with the topology of pointwise convergence. Just as in the proof of Proposition \ref{proposition:barycenter entropies}, now using Theorem \ref{th: large sample and catalytic}, we may easily prove that
\begin{equation}
\Gamma(X|Y)_P \leq \Gamma(X'|Y')_Q,
\;\; \forall \Gamma \in \mc H \;\; \implies \;\; \H(X|Y)_P \leq \H(X'|Y')_Q, \;\; \forall \H \;\text{in Definition~\ref{def: conditional entropy}} \,.
\end{equation}
Using this as a stepping stone, we may define a positive linear functional $G$ on the vector subspace of the set $C(\mc H)$ of continuous real functions on the compact Hausdorff $\mc H$ generated by the evaluation functions ${\rm ev}_{P}$, ${\rm ev}_P(\Gamma)=\Gamma(X|Y)_P$, through $G({\rm ev}_P)=\mb H(X|Y)_P$. Using Kantorovich's theorem we may extend $G$ into a positive linear functional $I:C(\mc H)\to\R$ and the Riesz-Markov-Kakutani theorem implies the existence of a finite positive measure $\nu:\mc B(\mc H)\to\R_+$ (Borel $\sigma$-algebra of the compact Hausdorff space $\mc H$) which is inner and outer regular such that
\begin{equation}
\mb H(X|Y)_P=\int_{\mc H}\Gamma(X|Y)_P\,\mathrm{d}\nu(\Gamma)\,.
\end{equation}
This measure is easily seen to be a probability measure.
\end{proof}


\bigskip


\subsection{Conditional entropies in the literature}
\label{sec: literature}
In this section, we review in detail several conditional entropies that have been introduced in the literature and show how each of these quantities arises as a special case of the general form established in Theorem~\ref{th:barycenter conditional entropies}.



The conditional Shannon entropy
\begin{equation}
    H(X|Y)_P=-\sum_{x \in \X,y\in\Y} P_{XY}(x,y)\log{\frac{P_{XY}(x,y)}{P_Y(y)}} \,,
\end{equation}
is equal to \(H_{0,\alpha}\) for \(\alpha = 1\).
The Hayashi conditional entropy introduced by Hayashi and Skorić et al.~\cite{hayashi2011exponential,vskoric2011sharp}
\begin{equation}
H^\text{H}_{\alpha}(X|Y)_P=\frac{1}{1-\alpha}\log{\Big(\sum\nolimits_{x\in\X,y\in\Y} P_Y(y) P_{X|Y}(x|y)^\alpha\Big)} \,,
\end{equation}
arises as a special case of $H_{t,\tau}$ when $\tau$ is the point measure at $\alpha$ and $t = 1 - \alpha$. The Arimoto Rényi conditional entropy~\cite{arimoto1977information}
\begin{align}
H^\text{A}_{\alpha}(X|Y)_P=\frac{\alpha}{1-\alpha}\log{\bigg(\sum_{y\in \Y} P_Y(y)\Big(\sum_{x \in \X} P_{X|Y}(x|y)^\alpha\Big)^\frac{1}{\alpha}\bigg)} \,,
\end{align}
is recovered similarly as a special case of $H_{t,\tau}$ when $\tau$ is the point measure at $\alpha$ and $t = \frac{1 - \alpha}{\alpha}$. A two-parameter family that interpolates between the Arimoto and Hayashi conditional entropies was introduced in~\cite{hayashi2016equivocations} and later investigated in detail in~\cite{rubboli2024quantum}
\begin{align}
\label{eq: H alpha beta}
H_{\alpha,\beta}(X|Y)_P=\frac{\alpha}{1-\alpha}\frac{1}{\beta}\log{\bigg(\sum_{y \in \Y} P_Y(y)\Big(\sum_{x \in \X} P_{X|Y}(x|y)^\alpha\Big)^\frac{\beta} {\alpha}\bigg)} \,.
\end{align}
This quantity is a special case of $H_{t,\tau}$ when $\tau$ is the point measure at $\alpha$ and $t = \frac{1 - \alpha}{\alpha} \beta$. More recently, this quantity has been studied in~\cite{li2025two}.


More exotic quantities have appeared in the literature. In~\cite{tan2018conditional}, the authors defined the conditional entropy
\begin{equation}
    H_{1+a|1+b}(X|Y)_P=-\frac{1+b}{a}\log{\bigg(\sum_{y \in \Y} P_Y(y)\Big(\sum_{x \in \X} P_{X|Y}(x|y)^{1+a}\Big)\Big(\sum_{x \in \X} P_{X|Y}(x|y)^{1+b}\Big)^{-\frac{a}{1+b}}\bigg)} \,.
\end{equation}
This corresponds to $H_{t,\tau}$ for $\tau=(1+b)\delta_{1+a}-b\delta_{1+b}$ and $t=-a/(1+b)$.
An alternative conditional entropy was introduced in~\cite{hayashi2016uniform}\footnote{Equation (14) of Ref.~\cite{hayashi2016uniform} appears to omit a minus sign.  We use the exponent $-\theta/(1+\theta')$, as the later arguments in that reference do.}
\begin{align}
    H_{1+\theta,1+\theta'}(X|Y)_P=&-\frac{1}{\theta}\log{\bigg(\sum_{y \in \Y} P_Y(y)\Big(\sum_{x \in \X} P_{X|Y}(x|y)^{1+\theta}\Big)\Big(\sum_{x \in \X} P_{X|Y}(x|y)^{1+\theta'}\Big)^{-\frac{\theta}{1+\theta'}}\bigg)} \\
    &\qquad + \frac{\theta'}{1+\theta'}H^{\text{A}}_{1+\theta'}(P_{XY}) \\
    =&\frac{1}{1+\theta'}H_{t_1,\tau_1}(X|Y)_P+\frac{\theta'}{1+\theta'}H_{t_2,\tau_2}(X|Y)_P\,,
\end{align}
Here, $\tau_1=(1+\theta')\delta_{1+\theta}-\theta'\delta_{1+\theta'}$, $t_1=-\theta/(1+\theta')$, $\tau_2$ is the point measure at $1+\theta'$, and $t_2=-\theta'/(1+\theta')$.

The conditional entropy
 \begin{align}
H_{0,\alpha}(P)=\sum_{y\in\mc Y} P_Y(y)H_\alpha(P_{X|Y=y})\,
 \end{align}
 was introduced by Cachin~\cite{cachin1997entropy} and later studied in~\cite{teixeira2012conditional}.

In~\cite{renner2005simple},
as a part of a study of conditional smooth R\'enyi entropies, the authors introduced the quantity
\begin{align}
    H^{\text{R}}_\alpha(X|Y) = \min _{y} H_\alpha(X|Y=y)
\end{align}
for $\alpha > 1$ and with $\min$ replaced by $\max$ for $\alpha \in (0,1)$. In the range $\alpha>1$, it coincides with $H_{-\infty,\tau}$ when $\tau$ is a point measure at $\alpha$.
For $\alpha \in (0,1)$, it coincides with the limit $t\rightarrow +\infty$ of $H_{t,\tau}$ in the case where $\tau$ is a point measure at $\alpha$. However, although this limit is true for all $\alpha \in (0,1)$, Theorem~\ref{th:barycenter conditional entropies} implies that these quantities satisfy all our axioms required for a conditional entropy  (Definition~\ref{def: conditional entropy}) only in the case $\alpha = 0$.
In this case, it recovers \(H_{+\infty,0}\).

Finally, the quantities $ H_{+\infty,0}$ and $H_{-\infty,\tau}$ for single-point measures have been studied as limiting cases of
$H_{\alpha,\beta}$~\eqref{eq: H alpha beta} in~\cite{li2025two} and~\cite{rubboli2024quantum}.





\section{Conclusion}

We characterized conditional entropy from four operational requirements:
invariance under relabeling and embedding, monotonicity under conditionally
mixing channels, additivity, and normalization.  Every such entropy is a
barycentre of the extremal families $H_{t,\tau}$ and their tropical and
derivation limits.  The preordered-semiring spectrum supplies the candidate
families, while Sections~4 and~5 determine exactly which parameters are
monotone.

The parameter proof treats arbitrary probability measures on the compactified
R\'enyi interval and does not assume finiteness of the singular moment at the
Shannon point.  In the exceptional negative branch, truncated localization
first forces the support restriction and only then proves the required
one-sided finiteness.  The complete analytic calculations and the
machine-checked statement correspondence are available in
Refs.~\cite{rubboli2026parametersFull,rubboli2026parametersLean},
respectively.




\section{Acknowledgements}
We thank Frits Verhagen for helpful discussions. R.R.\ was supported by the National Research Foundation, Singapore, and A*STAR under its CQT Bridging Grant and acknowledges financial support from the ERC grant GIFNEQ 101163938. E.H.\ and M.T.\ acknowledge support from the National Research Foundation Investigatorship Award (NRF-NRFI10-2024-0006) and from the National Research Foundation, Singapore through the National Quantum Office, hosted in A*STAR, under its Centre for Quantum Technologies Funding Initiative (S24Q2d0009).


\bibliographystyle{ultimate}
\bibliography{bibliography}


\newpage
\appendix

\section{Convexity and monotonicity}
\label{app: convexity and monotonicity}
In this appendix, we present several known equivalences between convexity and monotonicity under channels. These will be useful in the main text to characterize the parameter regimes for which the conditional entropy is monotone under conditionally mixing channels, and thus qualifies as a valid conditional entropy according to Definition~\ref{def: conditional entropy}. In particular, these equivalences allow us to reduce the question of monotonicity to that of convexity, which is often easier to analyze.

In Definition~\ref{def: conditional entropy}, we introduced conditional entropies as quantities satisfying a collection of operational axioms, motivated by the requirement that they quantify conditional uncertainty. A central axiom is monotonicity under conditionally mixing channels.
Within the mathematical framework of preordered semirings, monotone homomorphisms and derivations are functions designed to satisfy closely analogous structural properties (see Section~\ref{sec: background semiring}). In particular, in the semiring associated with conditional majorization, these functions correspond—up to an additive constant and a logarithmic factor—to (extremal) conditional entropies.
Consequently, understanding their monotonicity properties is essential. In the following, we show that determining whether these functions are monotone under conditionally mixing channels is equivalent to studying their convexity properties.


We begin by recalling several notions of convexity for functions defined on
(unnormalized) joint distributions~$\mathcal V$, which—up to embedding in a larger space (padding with zeros)—
can be identified with the elements of the semiring of conditional majorization
introduced in Section~\ref{sec: semiring of conditional entropies}.
 We say that a map $F:\V\to\mb R$ is convex if, for any $P,Q\in\V$ of the same size and $t\in[0,1]$, we have
\begin{equation}
F\big(tP+(1-t)Q\big)\leq tF(P)+(1-t)F(Q).
\end{equation}
We say that $F$ is concave if the above inequality is reversed. We say that $F$ is quasi-convex if, for $P,Q\in\V$ of the same size and $t\in(0,1)$,
\begin{equation}
F\big(tP+(1-t)Q\big)\leq\max\{F(P),F(Q)\}.
\end{equation}
We say that $F$ is strongly quasi-concave if the inequality above is reversed, i.e.,
\begin{equation}
\label{eq: def strongly quasi-conc}
    F\big(tP+(1-t)Q\big)\geq\max\{F(P),F(Q)\}.
\end{equation}

 In the following, we show that the required convexity properties depend on the output space of the homomorphism. As used in Sections~4 and~5 and detailed in Ref.~\cite{rubboli2026parametersFull}, the key monotonicity is monotonicity under maps acting on the conditioning subsystem of a joint distribution, which form the key subclass of conditionally mixing channels. In matrix form, for a positive joint distribution $P_{XY}$, such an action is represented as $P_{XY}\rightarrow P_{XY}D$, where $D$ is row stochastic.

\begin{proposition}\label{prop:convconc}
Let $S$ be the semiring of conditional majorization of Section~\ref{sec: semiring of conditional entropies}. Let
$\Phi: \cS \to \mathbb{K}$ be a candidate monotone homomorphism, and hence of the form in Proposition~\eqref{prop: monotone homomorphisms}. Here, $\mathbb{K} \in \{\mathbb{R}_+,  \mathbb{R}_+^{\mathrm{op}},\mathbb{TR}_+,\mathbb{TR}_+^{\mathrm{op}}\}$. Then, monotonicity with respect to maps acting on the conditioning system is equivalent to the following convexity properties, depending on the set $\mathbb{K}$:
\begin{itemize}[itemsep=0pt, topsep=5pt]
\item If $\mathbb{K} =\mathbb{R}_+$, $\Phi$ is monotone if and only if it is concave.
\item If $ \mathbb{K}= \mathbb{R}_+^{\mathrm{op}}$, $\Phi$ is monotone if and only if it is convex.
\item If $ \mathbb{K}=\mathbb{TR}_+$, $\Phi$ is monotone if and only if it is strongly quasi-concave.
\item If $ \mathbb{K}=\mathbb{TR}_+^{\mathrm{op}}$, $\Phi$ is monotone if and only if it is quasi-convex.
\end{itemize}
\end{proposition}

\begin{proof}
Let us first prove the `only if' part of the claim. The proof relies on associating the distributions $tP$ and $(1 - t)Q$ with distinct outcomes of the conditioning variable $Y$. Monotonicity under discarding outcomes of $Y$ then yields the desired inequality. Since monotone homomorphisms are invariant under embedding (see the discussion in
Section~\ref{sec: invariance embedding and shifting}), we may, in the following,
identify their evaluation on equivalence classes with their evaluation on
unnormalized joint distributions.


Explicitly, suppose that $P,Q\in\V$ and $t\in(0,1)$. We now have
\begin{equation}\label{eq:affinechain}
\left(\begin{array}{cc}
tP&0\\
0&(1-t)Q
\end{array}\right)\succeq\left(\begin{array}{cc}
tP&(1-t)Q\\
0&0
\end{array}\right)\succeq\left(\begin{array}{cc}
tP+(1-t)Q&0\\
0&0
\end{array}\right).
\end{equation}
The first implication follows by applying a permutation on the $X$ register conditioned on the second $Y$ outcome, while the second follows by discarding both outcomes and introducing a dummy outcome with zero probability.

In case $\mathbb{R}_+$, it now follows that
\begin{align}
\Phi\big(tP+(1-t)Q\big)&\geq\Phi\left[\left(\begin{array}{cc}
tP&0\\
0&(1-t)Q
\end{array}\right)\right]\\
&=\Phi(tP)+\Phi\big((1-t)Q\big)\\
&=\Phi(t)\Phi(P)+\Phi(1-t)\Phi(Q)\\
&=t\Phi(P)+(1-t)\Phi(Q).
\end{align}
Similarly, in case $\mathbb{R}_+^{\mathrm{op}}$, we get
\begin{align}
\Phi\big(tP+(1-t)Q\big)&\leq t\Phi(P)+(1-t)\Phi(Q) \,.
\end{align}
In case $\mathbb{TR}_+$, we get
\begin{align}
\Phi\big(tP+(1-t)Q\big)&\geq\max\{\Phi(P),\Phi(Q)\} \,.
\end{align}
In case $\mathbb{TR}^{\rm{op}}_+$, we have
\begin{align}
\Phi\big(tP+(1-t)Q\big)&\leq\Phi\left[\left(\begin{array}{cc}
tP&0\\
0&(1-t)Q
\end{array}\right)\right]\\
&=\max\{\Phi(tP),\Phi\big((1-t)Q\big)\}\\
&=\max\{\Phi(t)\Phi(P),\Phi(1-t)\Phi(Q)\}\\
&=\max\{\Phi(P),\Phi(Q)\}.
\end{align}

Next, we prove the “if” direction and assume that $\Phi$ is concave. Let $P \in \mathcal{V}$ and let $D$ be a row-stochastic matrix. We aim to show that $\Phi(PD) \geq \Phi(P)$. Let us first assume that $P=R=R_{X:YZ}$ (colon separating the primary $X$ from the conditioning $YZ$) and that $D$ is the map ignoring $Z$. %Just as in Lemma~\ref{lemma:shifting}, we see that
%\begin{equation}
%R\succeq\bigoplus_{z\in Z}R_Z(z)R_{XY|Z=z}\succeq R,
%\end{equation}
%so that the form given in \eqref{eq:FinalForm} tells us that
We have
\begin{align}
\Phi(R)&=\sum_{z\in\mc Z}R_Z(z)\Phi(R_{XY|Z=z})\leq\Phi\left(\sum_{z\in\mc Z}R_Z(z)R_{XY|Z=z}\right)=\Phi(R_{XY}),
\end{align}
where we have used the form of \eqref{eq:FinalForm} in the first equality and the assumed concavity of $\Phi$ in the inequality. This proves the first special case. Assume then the case $P_{XY}\succeq P_{XY}D=:Q_{XY'}$ where $D$ is a row-stochastic matrix which has only entries 0 or 1. This means that the matrix entries of $D$ are given by $D(y'|y)=\delta_{f(y),y'}$ for some function $f:Y\to Y'$. Let us define $R=R_{X:YY'}$ where $R(x,y,y')=P(x,y)$ whenever $f(y)=y'$ and $R(x,y,y')=0$ otherwise. It immediately follows that $P\approx R$ and $R_{XY'}=Q=PD$. Thus, using the above result, we have
\begin{align}
\Phi(PD)=\Phi(R_{XY'})\geq\Phi(R)=\Phi(P).
\end{align}
Finally, we prove the case of a general row-stochastic matrix $D$. Since $D$ can be expressed as a convex combination of row-stochastic matrices $D_k$ whose entries are either 0 or 1, there exist weights $t_k > 0$ with $\sum_k t_k = 1$
such that
\begin{equation}
    D = \sum_k t_k D_k.
\end{equation}
Our observations thus far, together with the concavity of $\Phi$ imply
\begin{align}\label{eq:apuconv}
\Phi(PD)&=\Phi\left(\sum_k t_kPD_k\right)\geq\sum_k t_k\Phi(PD_k)=\sum_k t_k\Phi(P)=\Phi(P).
\end{align}
The other cases are similar.
\end{proof}

A similar argument applies to derivations, which must be concave.

\begin{lemma}
\label{lem: conv derivations}
Let $S$ be the semiring of conditional majorization of Section~\ref{sec: semiring of conditional entropies}. Let $\Delta: \mathcal{S} \to \mathbb{R}$ be a candidate extremal monotone derivation at $\|\cdot\|$, and hence of the form in Proposition~\ref{prop: monotone derivations}.
Then, $\Delta$ is a monotone derivation if and only if it is concave.
\end{lemma}

\begin{proof}
By assumption, $\Delta:\cS\to\mb R$ is a derivation at $\|\cdot\|$ of the form in Proposition~\ref{prop: monotone derivations}. In particular, it satisfies $\Delta(x)=0$ for all $x\in\mb R_+$. Let now $P,Q\in\V$ be of the same size and $t\in[0,1]$. Using the same argument in the proof of Proposition~\ref{prop:convconc}, and in particular the relation in equation~\eqref{eq:affinechain}, we have
\begin{align}
\Delta\big(tP+(1-t)Q\big)&\geq\Delta(tP)+\Delta\big((1-t)Q\big)\\
&=\underbrace{\Delta(t)}_{=0}\|P\|+t\Delta(P)+\underbrace{\Delta(1-t)}_{=0}\|Q\|+(1-t)\Delta(Q)\\
&=t\Delta(P)+(1-t)\Delta(Q).
\end{align}

For the reverse implication, the monotonicity of $\Delta$ can be shown in exactly the same way as in the second part of the proof of Proposition~\ref{prop:convconc}.
\end{proof}

The next lemma shows that, for candidate homomorphisms, monotonicity under arbitrary maps acting on the conditioning system already implies monotonicity under the full class of conditionally mixing channels. This follows from the additional facts that the candidate homomorphisms are monotone under mixing channels and additive under direct sums. Together, these properties suffice to establish monotonicity under general conditionally mixing channels.
\begin{lemma}
\label{lem: extension to general channels}
Let $S$ be the semiring of conditional majorization of Section~\ref{sec: semiring of conditional entropies}.
Let $\Phi: \cS \to \mathbb{K}$ be a candidate monotone homomorphism, and hence of the form in Proposition~\ref{prop: monotone homomorphisms}. If $\Phi$ is monotone under row stochastic maps acting on the conditioning system, then it is monotone under conditionally mixing operations.
\end{lemma}

\begin{proof}
We present the proof for the case $\mathbb{K} = \mathbb{R}_+$, as the other cases follow analogously.

We first note that, as discussed in Appendix~\ref{app: convexity and monotonicity},
monotonicity under mixing channels is guaranteed by the explicit form of
the homomorphisms derived in Proposition~\ref{prop: monotone homomorphisms}, since
these are given by exponential expectations of R\'enyi entropies, which are
monotone under such maps.


To establish monotonicity under conditionally mixing channels, we introduce an auxiliary register that separates the individual terms appearing in the sum defining the channel, and then exploit the direct-sum property of the homomorphisms. The proof then follows by applying monotonicity under the action of doubly stochastic matrices and row-stochastic ones to each term separately.
Explicitly, let $P_{XY}\in\V$, $S^{(i)}$ be doubly stochastic matrices, and $D^{(i)}$ have non-negative entries ($i=1,\ldots,n$) with $D^{(1)}+\cdots+D^{(n)}$ row-stochastic. Let us denote the output of the conditionally mixing channel $Q_{X'Y'}=\sum_i S^{(i)}P_{XY}D^{(i)}$. Let us define new conditioning systems $\widetilde{Y}:=Y\times\{1,\ldots,n\}$ and $\widetilde{Y}':=Y'\times\{1,\ldots,n\}$.
Let $\pi_{Y'|\widetilde{Y}'}$ be the row-stochastic matrix which ignores the part $\{1,\ldots,n\}$ in $\widetilde{Y}'$. Let us also define the matrices
\begin{align}
P_{X\widetilde{Y}}&:=\big(P_{XY}D^{(1)}\,\cdots\,P_{XY}D^{(n)}\big)\\
Q_{X'\widetilde{Y}'}&:=\big(S^{(1)}P_{XY}D^{(1)}\,\cdots\,S^{(n)}P_{XY}D^{(n)}\big).
\end{align}
Because of the form of \eqref{eq:FinalForm}, %shifting Lemma in~\ref{lemma:shifting},
we have
\begin{equation}
\Phi(P_{X\widetilde{Y}})=\Phi\left(\bigoplus_{i=1}^n P_{XY}D^{(i)}\right),\quad\Phi(Q_{X'\widetilde{Y}'})=\Phi\left(\bigoplus_{i=1}^n S^{(i)}P_{XY}D^{(i)}\right).
\end{equation}
Putting all this together and noting that $(D^{(1)}\,\cdots D^{(n)})$ is row-stochastic and that $\bigoplus_{i=1}^n S^{(i)}$ is doubly stochastic, we have
\begin{align}
\Phi(P_{XY})&\leq\Phi\big[P_{XY}\big(D^{(1)}\,\cdots\,D^{(n)}\big)\big]\\
&=\Phi(P_{X\widetilde{Y}})=\Phi\left(\bigoplus_{i=1}^n P_{XY}D^{(i)}\right)\\
&\leq\Phi\left[\left(\bigoplus_{i=1}^n S^{(i)}\right)\left(\bigoplus_{i=1}^n P_{XY}D^{(i)}\right)\right]\\
&=\Phi\left(\bigoplus_{i=1}^n S^{(i)}P_{XY}D^{(i)}\right)\\
&=\Phi(Q_{X'\widetilde{Y}'})\\
&\leq\Phi(Q_{X'\widetilde{Y}'}\pi_{Y'|\widetilde{Y}'})\\
&=\Phi(Q_{X'Y'}),
\end{align}
showing the monotonicity in case $\mathbb{R}_+$. The other cases are essentially the same; the sums in \eqref{eq:apuconv} turn into maxima in the tropical cases $\mathbb{TR}_+$ and $\mathbb{TR}^{\mathrm{op}}_+$, and the inequalities are reversed in cases $\mathbb{R}^{\mathrm{op}}_+$ and $\mathbb{TR}^{\mathrm{op}}_+$.
\end{proof}
\begin{lemma}
\label{lem:extension-to-general-channels-derivations}
An extremal derivation at $\|\cdot\|$ that is monotone under row-stochastic
maps on the conditioning system is monotone under conditionally mixing
operations.
\end{lemma}
\begin{proof}
Repeat the proof of Lemma~\ref{lem: extension to general channels}, using
additivity of the derivation in place of additivity of the homomorphism.
\end{proof}

\end{document}
